% ROZDZIAL 11 — Wyjatki i bezpieczenstwo wyjatkowe w C++20
% Fragment rozdzialu do skryptu. Wymagane pakiety: listings, amsmath, array, hyperref.
% Uzycie: % ROZDZIAL 11 — Wyjatki i bezpieczenstwo wyjatkowe w C++20
% Fragment rozdzialu do skryptu. Wymagane pakiety: listings, amsmath, array, hyperref.
% Uzycie: % ROZDZIAL 11 — Wyjatki i bezpieczenstwo wyjatkowe w C++20
% Fragment rozdzialu do skryptu. Wymagane pakiety: listings, amsmath, array, hyperref.
% Uzycie: % ROZDZIAL 11 — Wyjatki i bezpieczenstwo wyjatkowe w C++20
% Fragment rozdzialu do skryptu. Wymagane pakiety: listings, amsmath, array, hyperref.
% Uzycie: \input{rozdzial_11_wyjatki_bezpieczenstwo_cpp.tex}

\chapter{Wyjątki i bezpieczeństwo wyjątkowe w C++}
\label{chap:cpp-exceptions}

\section{Po co C++ ma wyjątki?}

Mechanizm wyjątków oddziela kod realizujący zwykły przebieg algorytmu
od kodu reagującego na niepowodzenia. Jeżeli funkcja nie może spełnić
swojego kontraktu, może zgłosić wyjątek poleceniem \texttt{throw}.
Wówczas wykonanie jest przekazywane do odpowiedniego bloku
\texttt{catch}; po drodze niszczone są obiekty automatyczne, których
czas życia dobiegł końca. To ostatnie zjawisko nazywamy
\emph{odwijaniem stosu} (\emph{stack unwinding}).

Najprostszy schemat wygląda tak: \texttt{throw} zgłasza problem,
a pasujący \texttt{catch} go odbiera. Kod po \texttt{throw} w tym samym
bloku nie zostanie wykonany:
\begin{lstlisting}
try {
    throw std::runtime_error("bad input");
} catch (const std::runtime_error& error) {
    std::cerr << error.what() << '\n';
}
\end{lstlisting}
W samodzielnym programie ten fragment wymaga nagłówków
\texttt{<iostream>} i \texttt{<stdexcept>}.

Z perspektywy projektowania API warto odróżniać:
\begin{itemize}
\item \emph{zwykły wynik} (w tym informację, że elementu nie znaleziono),
      dla którego często lepszy jest \texttt{std::optional};
\item \emph{błąd operacji, który wywołujący powinien świadomie obsłużyć},
      reprezentowany np. przez kod błędu lub, od C++23, \texttt{std::expected};
\item \emph{niepowodzenie uniemożliwiające dotrzymanie kontraktu},
      dla którego wyjątek bywa naturalnym sposobem propagacji.
\end{itemize}
Nie ma uniwersalnej reguły nakazującej używanie wyjątków wszędzie.
Znaczenie mają kontrakty funkcji, konwencje projektu, ograniczenia
środowiska oraz wymagania dotyczące opóźnień wykonania.

\section{Składnia \texttt{try}, \texttt{throw} i \texttt{catch}}

\begin{center}\small
\begin{tabularx}{\textwidth}{@{}p{0.32\textwidth}X@{}}
\toprule
Konstrukcja & Za co odpowiada \\
\midrule
\code{try} & Obejmuje blok, z którego wyjątki mogą zostać przechwycone przez następujące bloki \code{catch}. \\
\code{throw value} & Tworzy obiekt wyjątku i rozpoczyna poszukiwanie pasującej obsługi. \\
\texttt{catch(const E\& e)} & Przechwytuje wyjątek zgodnego typu bez kopiowania i bez obcinania części pochodnej. \\
\code{catch(...)} & Przechwytuje dowolny wyjątek, także nieznanego typu. \\
\code{throw;} & Wewnątrz obsługi ponownie zgłasza bieżący wyjątek, zachowując jego typ dynamiczny. \\
\code{std::exception::what()} & Zwraca tekstowy opis błędu; służy do diagnostyki, nie do sterowania logiką. \\
\bottomrule
\end{tabularx}
\end{center}
\label{sec:cpp11-basics}

\begin{lstlisting}
#include <iostream>
#include <stdexcept>

double divide(double a, double b) {
    if (b == 0.0)
        throw std::invalid_argument("zero denominator");
    return a / b;
}

int main() {
    try {
        std::cout << divide(10.0, 2.0) << '\n';
        std::cout << divide(10.0, 0.0) << '\n';
        std::cout << "not executed\n";
    } catch (const std::invalid_argument& e) {
        std::cerr << "invalid input: " << e.what() << '\n';
    }
}
\end{lstlisting}

Instrukcja \texttt{throw E\{...\}} tworzy obiekt wyjątku. Obsługa
szuka pasującej klauzuli \texttt{catch}; jeśli wyjątek nie zostanie
obsłużony, program kończy działanie poprzez \texttt{std::terminate}.
Nie należy zakładać, że nieobsłużony wyjątek zawsze spowoduje pełne
odwinięcie stosu: szczegóły zakończenia programu mogą zależeć od
implementacji.

\subsection{Przechwytywanie przez referencję i kolejność klauzul}

Przyjmujemy zwykle \texttt{const E\&}: pozwala to uniknąć niepotrzebnego
kopiowania i \emph{object slicing} przy przekazywaniu obiektu typu
pochodnego jako typu bazowego przez wartość.

\begin{lstlisting}
#include <exception>
#include <iostream>
#include <stdexcept>

void handle() {
    try {
        throw std::out_of_range("bad index");
    } catch (const std::out_of_range& e) {
        std::cout << "out_of_range: " << e.what() << '\n';
    } catch (const std::exception& e) {
        std::cout << "other std::exception: " << e.what() << '\n';
    } catch (...) {
        std::cout << "unknown exception\n";
    }
}
\end{lstlisting}

Klauzule są sprawdzane w kolejności występowania: najpierw umieszcza
się typy bardziej szczegółowe, a na końcu bazowe. Wyrażenie
\texttt{catch (...)} przechwytuje wyjątek dowolnego typu, ale nie
udostępnia bezpośrednio jego obiektu.

\paragraph{Ponowne zgłoszenie.}
\texttt{throw;} wewnątrz aktywnej obsługi ponownie zgłasza bieżący
wyjątek z zachowaniem jego dynamicznego typu. Natomiast
\texttt{throw e;} zgłasza nowe wyrażenie; w zależności od statycznego
typu \texttt{e} może prowadzić do kopiowania i obcięcia do klasy
bazowej.

\begin{lstlisting}
void log_and_rethrow() {
    try {
        throw std::runtime_error("disk error");
    } catch (const std::exception& e) {
        std::cerr << e.what() << '\n';
        throw;                  // keeps original exception object
        // throw e;              // may slice to std::exception
    }
}
\end{lstlisting}

\section{Odwijanie stosu i RAII}
\label{sec:cpp11-unwinding}

Najważniejsza reguła praktyczna brzmi: \emph{zwalnianie zasobu powinno
wynikać z destruktora obiektu zarządzającego zasobem}. Stanowi to
rozwinięcie idei RAII poznanej w rozdziale~\ref{chap:wskazniki-cpp}.

\begin{lstlisting}
#include <iostream>
#include <stdexcept>
#include <string>
#include <utility>

struct Trace {
    std::string name;
    explicit Trace(std::string n) : name(std::move(n)) {
        std::cout << "create " << name << '\n';
    }
    ~Trace() noexcept {
        std::cout << "destroy " << name << '\n';
    }
};

void third() {
    Trace c{"C"};
    throw std::runtime_error("failure");
}
void second() {
    Trace b{"B"};
    third();
}
void first() {
    Trace a{"A"};
    second();
}
// try { first(); } catch (const std::exception&) {}
// Destruction order: C, B, A.
\end{lstlisting}

Wyjątek powoduje zakończenie odpowiednich zakresów, a obiekty
z lokalnym automatycznym czasem przechowywania są niszczone w odwrotnej
kolejności ich skonstruowania. Jeżeli konstruktor zgłosi wyjątek,
destruktor \emph{całego, nieutworzonego obiektu} nie jest wywoływany,
ale niszczone są już w pełni skonstruowane podobiekty bazowe i składowe.

\subsection{Zarządzanie plikiem przy pomocy RAII}

\begin{lstlisting}
#include <cstdio>
#include <memory>
#include <stdexcept>

struct FileCloser {
    void operator()(std::FILE* file) const noexcept {
        if (file) std::fclose(file);
    }
};
using File = std::unique_ptr<std::FILE, FileCloser>;

File open_file(const char* name) {
    File file{std::fopen(name, "rb")};
    if (!file) throw std::runtime_error("cannot open file");
    return file;
}
\end{lstlisting}

Nawet gdy po otwarciu pliku dalsze operacje zgłoszą wyjątek,
\texttt{unique\_ptr} wykona deleter. Zauważmy, że
\texttt{std::fclose} może sygnalizować błąd przy zamykaniu; prosty
noexcept-deleter ignoruje ten kod błędu. Jeżeli informacja o błędzie
zamknięcia jest istotna, należy zapewnić osobną, jawną operację
finalizacji, której wynik można obsłużyć \emph{przed} destrukcją.

\section{Gwarancje bezpieczeństwa wyjątkowego}

\begin{center}\small
\begin{tabularx}{\textwidth}{@{}p{0.32\textwidth}X@{}}
\toprule
Gwarancja lub narzędzie & Skutek w razie wyjątku \\
\midrule
Brak wycieku & Zasoby pozostają poprawnie zwalniane, zwykle przez RAII. \\
Gwarancja podstawowa & Obiekty pozostają poprawne, choć ich stan może się zmienić. \\
Gwarancja silna & Operacja kończy się sukcesem albo pozostawia obserwowalny stan bez zmian. \\
Gwarancja braku wyjątków & Operacja nie pozwala wyjątkowi opuścić funkcji. \\
\code{noexcept} & Deklaruje gwarancję braku wyjątku; wyjątek opuszczający funkcję kończy program przez \code{std::terminate}. \\
\bottomrule
\end{tabularx}
\end{center}
\label{sec:cpp11-guarantees}

Bezpieczeństwo wyjątkowe nie oznacza tylko braku wycieków pamięci.
Załóżmy, że lista zawiera \texttt{[1, 2]} i dopisanie \texttt{3}
zgłosi wyjątek. Gwarancja podstawowa oznacza, że lista nadal nadaje
się do użycia, choć może mieć zmienioną zawartość. Gwarancja silna
oznacza, że nadal zawiera dokładnie \texttt{[1, 2]}. Gwarancja braku
wyjątków mówi tylko, że wyjątek nie opuści funkcji; jeśli operacja
ma dodatkowo zawsze osiągnąć cel, mówimy o gwarancji \emph{no-fail}.

Silna gwarancja nie oznacza automatycznie cofnięcia efektów
zewnętrznych, np. wysłanego pakietu czy wpisu w bazie danych.
Zakres transakcyjności trzeba sprecyzować w kontrakcie.

\subsection{Silna gwarancja dzięki operacji „przygotuj, a potem zatwierdź”}

\begin{lstlisting}
#include <utility>
#include <vector>

class Table {
    std::vector<int> data_;
public:
    explicit Table(std::vector<int> values)
        : data_(std::move(values)) {}

    void append_many(const std::vector<int>& more) {
        auto candidate = data_;          // may throw
        candidate.insert(candidate.end(),
                         more.begin(), more.end()); // may throw
        data_.swap(candidate);           // noexcept for std::vector<int>
    }

    const std::vector<int>& data() const noexcept { return data_; }
};
\end{lstlisting}

Jeżeli kopiowanie lub \texttt{insert} zawiedzie, \texttt{data\_}
pozostaje nietknięte. Dopiero bezwyjątkowe \texttt{swap} zatwierdza
przygotowaną wersję. Kosztem jest dodatkowa pamięć i kopiowanie;
silna gwarancja bywa więc świadomym wyborem projektowym, a nie darmową
optymalizacją.

\paragraph{Uwaga na aliasing.}
W przykładzie argument \texttt{more} może nawet odnosić się do
\texttt{data\_} obiektu \texttt{Table}: kopia \texttt{candidate}
powstaje przed modyfikacją oryginału. W kodzie operującym bezpośrednio
na tym samym kontenerze trzeba dodatkowo analizować ważność iteratorów
i referencji.

\subsection{Konstruktor i niezmienniki klasy}

Dobrze napisany konstruktor tworzy w pełni poprawny obiekt albo
zgłasza wyjątek. Bezpiecznie zarządzane składowe automatycznie
posprzątają po częściowo wykonanej inicjalizacji.

\begin{lstlisting}
#include <memory>
#include <stdexcept>

class Session {
    std::unique_ptr<int[]> buffer_;
    std::size_t size_;
public:
    explicit Session(std::size_t n)
        : buffer_(std::make_unique<int[]>(n)), size_(n) {
        if (n == 0) throw std::invalid_argument("empty session");
    }
    std::size_t size() const noexcept { return size_; }
};
\end{lstlisting}

Nie należy zastępować RAII ręcznie pisanymi \texttt{catch}, w których
sprząta się kolejne zasoby. Takie procedury są podatne na błędy,
zwłaszcza gdy konstruktor zdobywa wiele zasobów po kolei.

\section{\texttt{noexcept}: deklarowanie i sprawdzanie}
\label{sec:cpp11-noexcept}

\texttt{noexcept} w specyfikacji funkcji oznacza, że wyjątek nie może
opuścić jej ciała. Naruszenie tego wymagania kończy się wywołaniem
\texttt{std::terminate} (nie jest to „ciche zignorowanie” błędu).

\begin{lstlisting}
void close_handle() noexcept {
    // cleanup that must not let exceptions escape
}

void may_fail();
static_assert(noexcept(close_handle()));
static_assert(!noexcept(may_fail()));
\end{lstlisting}

W drugim użyciu \texttt{noexcept(expr)} jest \emph{operatorem}
zwracającym wartość logiczną w czasie kompilacji. Nie uruchamia on
samego wyrażenia.

\subsection{Warunkowe \texttt{noexcept}}

\begin{lstlisting}
#include <type_traits>
#include <utility>

template<class T>
void swap_values(T& a, T& b)
    noexcept(std::is_nothrow_move_constructible_v<T> &&
             std::is_nothrow_move_assignable_v<T>) {
    T tmp(std::move(a));
    a = std::move(b);
    b = std::move(tmp);
}
\end{lstlisting}

Jest to przykład \emph{warunkowej specyfikacji}. Nie oznacza,
że przedstawiony algorytm daje silną gwarancję, jeżeli przenoszenie
rzuca wyjątki; dla takich typów stan \texttt{a} lub \texttt{b} może
się zmienić przed niepowodzeniem. W kodzie generycznym trzeba
analizować nie tylko składnię, ale też gwarancje operacji typu \texttt{T}.

\subsection{Dlaczego \texttt{noexcept} jest ważne dla \texttt{vector}?}

Przy powiększaniu pojemności kontener może potrzebować przenieść
lub skopiować stare elementy do nowej pamięci. Dla typów kopiowalnych
konstrukcja kopiująca bywa bezpieczniejsza, jeśli konstruktor
przenoszący może zgłosić wyjątek.

\begin{lstlisting}
#include <type_traits>
#include <utility>

struct Item {
    Item() = default;
    Item(const Item&) {}              // copying available
    Item(Item&&) noexcept(false) {}   // move may throw
};

static_assert(std::is_copy_constructible_v<Item>);
static_assert(!std::is_nothrow_move_constructible_v<Item>);

// std::move_if_noexcept(x) usually gives const Item& here.
\end{lstlisting}

\texttt{std::move\_if\_noexcept} zwraca \texttt{T\&\&}, jeżeli
przenoszenie jest zadeklarowane jako bezwyjątkowe albo typ nie jest
kopiowalny; w przeciwnym razie zwraca \texttt{const T\&}.
Zasada ta pomaga kontenerom utrzymywać silną gwarancję w wielu,
ale nie we wszystkich przypadkach. Szczegóły są zależne od
konkretnej operacji i własności typu elementów.

\subsection{Destruktory i \texttt{std::terminate}}

Domyślnie destruktor jest niewyrzucający, o ile destruktory jego
składowych i klas bazowych nie wymagają innej specyfikacji.
Szczególnie podczas odwijania stosu wypuszczenie kolejnego wyjątku
z destruktora może spowodować \texttt{std::terminate}.

\begin{lstlisting}
struct SafeCleanup {
    ~SafeCleanup() noexcept {
        try {
            risky_cleanup();
        } catch (...) {
            // Log, record an error, or use another no-throw fallback.
            // Never let it escape this destructor.
        }
    }
    static void risky_cleanup();
};
\end{lstlisting}

Nie oznacza to, że każdą awarię należy ukryć. Jeśli operacja
finalizacji może się nie udać i użytkownik musi znać wynik, projektujmy
jawną metodę \texttt{close()} lub \texttt{commit()}, a destruktor
niech pełni rolę bezpiecznej ścieżki awaryjnej.

\section{Własne klasy wyjątków i hierarchia standardowa}
\label{sec:cpp11-types}

W nagłówkach \texttt{<stdexcept>}, \texttt{<exception>},
\texttt{<new>} i \texttt{<system\_error>} znajdziemy wiele gotowych
typów. W typowych zastosowaniach najczęściej spotkamy
\texttt{std::invalid\_argument}, \texttt{std::out\_of\_range},
\texttt{std::runtime\_error}, \texttt{std::bad\_alloc} oraz
\texttt{std::system\_error}.

\begin{lstlisting}
#include <cstddef>
#include <stdexcept>
#include <string>
#include <utility>

class ParseError : public std::runtime_error {
    std::size_t position_;
public:
    ParseError(std::size_t pos, std::string message)
        : std::runtime_error(std::move(message)),
          position_(pos) {}

    std::size_t position() const noexcept { return position_; }
};

int parse_digit(char c, std::size_t pos) {
    if (c < '0' || c > '9')
        throw ParseError(pos, "expected decimal digit");
    return c - '0';
}
\end{lstlisting}

Własna klasa może przechowywać dodatkowe informacje, np. identyfikator
operacji, pozycję w tekście lub kod błędu. Nie należy umieszczać tam
niepotrzebnie ogromnych danych. Metoda \texttt{what()} klasy
\texttt{std::exception} jest wirtualna i \texttt{noexcept}; jeśli
ją przesłaniamy, musimy zachować zgodny kontrakt i poprawny czas życia
zwracanego łańcucha.

\paragraph{Błędy systemowe.}
\texttt{std::system\_error} umożliwia przenoszenie kodu błędu
\texttt{std::error\_code} wraz z kategorią, np. błędu systemowego,
bez uzależniania całego API od przypadkowego tekstu \texttt{what()}.

\section{Wyjątki w funkcjach szablonowych i neutralność wyjątkowa}

Komponent generyczny powinien zwykle zachować niezmienniki i pozwolić
wyjątkom operacji wywołanych na typie \texttt{T} przejść do klienta,
zamiast pochłaniać je przez \texttt{catch (...)}.

\begin{lstlisting}
#include <algorithm>
#include <vector>

template<class T, class Compare>
void sort_elements(std::vector<T>& xs, Compare comp) {
    std::sort(xs.begin(), xs.end(), comp);
    // An exception from comp may propagate to the caller.
}
\end{lstlisting}

Należy jednak odróżnić \emph{neutralność wyjątkową} od silnej
gwarancji. Jeśli komparator zgłosi wyjątek w trakcie sortowania,
kontener pozostaje poprawnym obiektem, ale jego kolejność nie musi
odpowiadać stanowi sprzed wywołania.

\section{Przenoszenie wyjątków między funkcjami i wątkami}
\label{sec:cpp11-exception-ptr}

Ta i dwie kolejne sekcje pokazują zastosowania bardziej zaawansowane.
Do podstaw wystarczą \texttt{try}/\texttt{catch}, RAII oraz gwarancje
stanu po niepowodzeniu.

Typ \texttt{std::exception\_ptr} pozwala zachować wyjątek poza blokiem
\texttt{catch}, a później zgłosić go ponownie. Jest przydatny
szczególnie w programowaniu współbieżnym.

\begin{lstlisting}
#include <exception>
#include <iostream>
#include <stdexcept>
#include <string>
#include <thread>

std::exception_ptr error;
std::thread worker([&] {
    try {
        throw std::runtime_error("worker failed");
    } catch (...) {
        error = std::current_exception();
    }
});
worker.join();

if (error) {
    std::rethrow_exception(error);
}
\end{lstlisting}

Zmienna \texttt{error} nie jest tutaj współdzielona jednocześnie przez
dwa aktywne dostępy: zapis zachodzi w wątku roboczym, a odczyt
dopiero po \texttt{join()}, które zapewnia wymaganą synchronizację.
Nie wolno dopuścić, aby wyjątek opuścił funkcję startową zwykłego
\texttt{std::thread}: typowym skutkiem jest \texttt{std::terminate}.

\subsection{\texttt{std::future} przechowuje wyjątek}

Wspomniane w rozdziale o współbieżności \texttt{std::async} i
\texttt{std::future} ułatwiają propagowanie niepowodzenia do
odbiorcy wyniku.

\begin{lstlisting}
#include <future>
#include <stdexcept>

auto fut = std::async(std::launch::async, []() -> int {
    throw std::runtime_error("calculation failed");
});
try {
    (void)fut.get();  // rethrows the stored exception
} catch (const std::runtime_error&) {
    // report the failure in the receiving thread
}
\end{lstlisting}

Przez \texttt{get()} wynik (lub wyjątek) trafia do wątku odbiorcy.
W \texttt{std::future} wynik można pobrać tylko raz; inne własności
tego mechanizmu omówiliśmy w rozdziale o asynchroniczności.

\section{Wyjątki zagnieżdżone i dodatkowy kontekst}
\label{sec:cpp11-nested}

Gdy funkcja wyższego poziomu rozumie więcej niż warstwa niższa, może
wzbogacić informację o błędzie i zachować oryginalną przyczynę.

\begin{lstlisting}
#include <exception>
#include <stdexcept>

void load_config() {
    try {
        throw std::runtime_error("invalid token");
    } catch (...) {
        std::throw_with_nested(
            std::runtime_error("cannot load configuration"));
    }
}

void print_nested(const std::exception& e, int depth = 0) {
    std::cerr << std::string(depth * 2, ' ') << e.what() << '\n';
    try {
        std::rethrow_if_nested(e);
    } catch (const std::exception& inner) {
        print_nested(inner, depth + 1);
    } catch (...) {
        std::cerr << "nested non-standard exception\n";
    }
}
\end{lstlisting}

Takie opakowywanie wyjątków nie jest tym
samym co zgłoszenie \texttt{throw e;}: nowy wyjątek przechowuje
kontekst, a przyczynę można odzyskać przez \texttt{rethrow\_if\_nested}.

\section{\texttt{std::uncaught\_exceptions} i strażnicy zakresu}

Funkcja \texttt{std::uncaught\_exceptions()} (C++17) zwraca liczbę
aktualnie nieobsługiwanych obiektów wyjątków w bieżącym wątku.
Można jej użyć do wykrycia, czy wyjście z zakresu następuje wskutek
odwijania stosu.

\begin{lstlisting}
#include <exception>

class FailureGuard {
    int initial_ = std::uncaught_exceptions();
    int& failure_count_;
public:
    explicit FailureGuard(int& count) noexcept : failure_count_(count) {}
    FailureGuard(const FailureGuard&) = delete;
    FailureGuard& operator=(const FailureGuard&) = delete;

    ~FailureGuard() noexcept {
        if (std::uncaught_exceptions() > initial_) {
            ++failure_count_; // no-throw bookkeeping
        }
    }
};
\end{lstlisting}

Porównanie z \emph{początkową liczbą}, a nie tylko z zerem, jest
ważne przy zagnieżdżonych wyjątkach. Strażnik nie daje jednak
pełnego, transakcyjnego mechanizmu cofania zmian. W projektach
korzystających z C++20 można napisać własny scope guard, ale trzeba
uważnie określić, czy przechowuje funkcję przez kopię, przez
przeniesienie czy przez referencję i czy jej wywołanie może rzucać.

\section{Wyjątki a alternatywne modele obsługi błędów}
\label{sec:cpp11-alternatives}

\subsection{\texttt{std::optional}: brak wyniku nie musi być błędem}

\begin{lstlisting}
#include <optional>
#include <string_view>

std::optional<std::size_t> find_colon(std::string_view s) {
    const auto i = s.find(':');
    if (i == std::string_view::npos) return std::nullopt;
    return i;
}
\end{lstlisting}

Brak dwukropka może być normalnym scenariuszem. Zgłaszanie wyjątku
przy każdym nieudanym wyszukiwaniu bywa niepotrzebne.

\subsection{\texttt{std::expected<T,E>} — C++23, nie C++20}

\texttt{std::expected} reprezentuje \emph{albo wartość typu T, albo
błąd typu E}. Został ustandaryzowany w C++23 i nie należy zakładać,
że nagłówek \texttt{<expected>} jest dostępny w kompilacji C++20.

\begin{lstlisting}
// This example requires C++23 and <expected>.
#include <expected>
#include <string>
#include <string_view>

std::expected<int, std::string>
parse_positive(std::string_view text) {
    if (text == "1") return 1;
    if (text == "2") return 2;
    return std::unexpected("expected 1 or 2");
}

void use_expected() {
    auto result = parse_positive("3");
    if (!result) {
        const std::string& error = result.error();
        (void)error;
    }
}
\end{lstlisting}

Dzięki temu błędy oczekiwane przez API są jawne w sygnaturze.
\texttt{std::expected} nie eliminuje wszystkich wyjątków: konstrukcja
zawartego obiektu czy inne operacje nadal mogą się nie udać.
W C++20 podobny kontrakt można zrealizować własnym typem
\texttt{Result}, \texttt{std::variant<T,E>} lub odpowiednią biblioteką.

\subsection{Kody błędów, wyjątki i wydajność}

W konwencjonalnych implementacjach C++ obsługa sukcesu na ścieżce,
na której żaden wyjątek nie wystąpił, może być bardzo tania, ale samo
zgłoszenie i obsłużenie wyjątku zwykle kosztuje więcej niż zwrócenie
prostego kodu błędu. Nie należy jednak zamieniać każdego \texttt{throw}
na kod zwrotny tylko na podstawie intuicji: warto mierzyć rzeczywistą
częstotliwość błędów, profilować aplikację i uwzględniać prostotę
kontraktów.

\section{Typowe błędy i dobre praktyki}
\label{sec:cpp11-guidelines}

\begin{enumerate}
\item \textbf{Nie rzucaj wskaźnikami do lokalnych obiektów.}
      Preferuj wyjątki będące obiektami, zwykle pochodnymi od
      \texttt{std::exception}.
\item \textbf{Przechwytuj wyjątki polimorficzne przez
      \texttt{const\&}.} Unikaj obcinania obiektu (\emph{slicing}).
\item \textbf{Do ponownego zgłoszenia używaj \texttt{throw;}.}
      Nie zastępuj go bez namysłu \texttt{throw e;}.
\item \textbf{Stosuj RAII.} Surowe wywołania \texttt{new}/\texttt{delete}
      rozdzielone operacjami rzucającymi wyjątki to częsta przyczyna
      wycieków pamięci.
\item \textbf{Nie wypuszczaj wyjątków z destruktorów.}
      Błędy finalizacji zgłaszaj jawnie w odpowiednio zaprojektowanym API.
\item \textbf{Nie oznaczaj funkcji \texttt{noexcept} bez analizy.}
      To obietnica z konsekwencją w postaci \texttt{std::terminate}.
\item \textbf{Nie ukrywaj błędów w pustym \texttt{catch (...)}.}
      Obsłuż wyjątek, zapisz informację diagnostyczną albo przekaż go wyżej.
\item \textbf{Określ gwarancję niepowodzenia.} Dla operacji mutujących
      dane rozważ silną gwarancję i strategię „przygotuj–zatwierdź”.
\item \textbf{Złap wyjątek w funkcji wątku.}
      Nie pozwól mu opuścić funkcji startowej \texttt{std::thread}.
      Do przekazania błędu użyj \texttt{std::exception\_ptr},
      \texttt{promise/future} lub API wyższego poziomu.
\item \textbf{Wybierz właściwy model błędu.} Brak zwykłego wyniku,
      przewidywalna porażka operacji i naruszenie kontraktu to różne
      sytuacje.
\end{enumerate}

\section{Przykład końcowy: transakcyjna aktualizacja danych}

Poniższy program ilustruje współpracę wyjątku, RAII i silnej gwarancji.
Załóżmy, że aktualizacja dodaje wartości do rejestru, ale w razie błędu
rejestr ma pozostać bez zmian.

\begin{lstlisting}
#include <iostream>
#include <stdexcept>
#include <utility>
#include <vector>

class Register {
    std::vector<int> values_;
public:
    Register() : values_{1, 2, 3} {}

    void add_checked(const std::vector<int>& numbers) {
        std::vector<int> proposed = values_;
        for (int x : numbers) {
            if (x < 0) {
                throw std::invalid_argument("negative input");
            }
            proposed.push_back(x);
        }
        values_.swap(proposed);
    }

    const std::vector<int>& values() const noexcept {
        return values_;
    }
};

int main() {
    Register r;
    try {
        r.add_checked({4, -5, 6});
    } catch (const std::invalid_argument& e) {
        std::cerr << e.what() << '\n';
    }
    // Still {1, 2, 3}; no partial update.
    for (int x : r.values()) std::cout << x << ' ';
}
\end{lstlisting}

Najpierw powstaje kopia danych. Sprawdzanie i alokowanie nowej pamięci
mogą zgłaszać wyjątki, lecz działają na kopii. Po pomyślnym wykonaniu
następuje zatwierdzenie przez \texttt{swap}. Jest to wzorzec, który
można stosować również w rozbudowanych klasach, przy świadomym
uwzględnieniu kosztów i własności ich składowych.

\section{Podsumowanie rozdziału}

W praktyce mechanizm wyjątków jest najbardziej użyteczny wtedy,
gdy łączy się go z dobrym modelem własności zasobów i z precyzyjnie
określonymi gwarancjami operacji. \texttt{throw}, \texttt{catch}
i \texttt{noexcept} to tylko składnia; o poprawności projektu
decydują RAII, niezmienniki obiektów, semantyka transakcyjna oraz
świadome oddzielanie błędów wyjątkowych od zwykłych wariantów wyniku.



\chapter{Wyjątki i bezpieczeństwo wyjątkowe w C++}
\label{chap:cpp-exceptions}

\section{Po co C++ ma wyjątki?}

Mechanizm wyjątków oddziela kod realizujący zwykły przebieg algorytmu
od kodu reagującego na niepowodzenia. Jeżeli funkcja nie może spełnić
swojego kontraktu, może zgłosić wyjątek poleceniem \texttt{throw}.
Wówczas wykonanie jest przekazywane do odpowiedniego bloku
\texttt{catch}; po drodze niszczone są obiekty automatyczne, których
czas życia dobiegł końca. To ostatnie zjawisko nazywamy
\emph{odwijaniem stosu} (\emph{stack unwinding}).

Najprostszy schemat wygląda tak: \texttt{throw} zgłasza problem,
a pasujący \texttt{catch} go odbiera. Kod po \texttt{throw} w tym samym
bloku nie zostanie wykonany:
\begin{lstlisting}
try {
    throw std::runtime_error("bad input");
} catch (const std::runtime_error& error) {
    std::cerr << error.what() << '\n';
}
\end{lstlisting}
W samodzielnym programie ten fragment wymaga nagłówków
\texttt{<iostream>} i \texttt{<stdexcept>}.

Z perspektywy projektowania API warto odróżniać:
\begin{itemize}
\item \emph{zwykły wynik} (w tym informację, że elementu nie znaleziono),
      dla którego często lepszy jest \texttt{std::optional};
\item \emph{błąd operacji, który wywołujący powinien świadomie obsłużyć},
      reprezentowany np. przez kod błędu lub, od C++23, \texttt{std::expected};
\item \emph{niepowodzenie uniemożliwiające dotrzymanie kontraktu},
      dla którego wyjątek bywa naturalnym sposobem propagacji.
\end{itemize}
Nie ma uniwersalnej reguły nakazującej używanie wyjątków wszędzie.
Znaczenie mają kontrakty funkcji, konwencje projektu, ograniczenia
środowiska oraz wymagania dotyczące opóźnień wykonania.

\section{Składnia \texttt{try}, \texttt{throw} i \texttt{catch}}

\begin{center}\small
\begin{tabularx}{\textwidth}{@{}p{0.32\textwidth}X@{}}
\toprule
Konstrukcja & Za co odpowiada \\
\midrule
\code{try} & Obejmuje blok, z którego wyjątki mogą zostać przechwycone przez następujące bloki \code{catch}. \\
\code{throw value} & Tworzy obiekt wyjątku i rozpoczyna poszukiwanie pasującej obsługi. \\
\texttt{catch(const E\& e)} & Przechwytuje wyjątek zgodnego typu bez kopiowania i bez obcinania części pochodnej. \\
\code{catch(...)} & Przechwytuje dowolny wyjątek, także nieznanego typu. \\
\code{throw;} & Wewnątrz obsługi ponownie zgłasza bieżący wyjątek, zachowując jego typ dynamiczny. \\
\code{std::exception::what()} & Zwraca tekstowy opis błędu; służy do diagnostyki, nie do sterowania logiką. \\
\bottomrule
\end{tabularx}
\end{center}
\label{sec:cpp11-basics}

\begin{lstlisting}
#include <iostream>
#include <stdexcept>

double divide(double a, double b) {
    if (b == 0.0)
        throw std::invalid_argument("zero denominator");
    return a / b;
}

int main() {
    try {
        std::cout << divide(10.0, 2.0) << '\n';
        std::cout << divide(10.0, 0.0) << '\n';
        std::cout << "not executed\n";
    } catch (const std::invalid_argument& e) {
        std::cerr << "invalid input: " << e.what() << '\n';
    }
}
\end{lstlisting}

Instrukcja \texttt{throw E\{...\}} tworzy obiekt wyjątku. Obsługa
szuka pasującej klauzuli \texttt{catch}; jeśli wyjątek nie zostanie
obsłużony, program kończy działanie poprzez \texttt{std::terminate}.
Nie należy zakładać, że nieobsłużony wyjątek zawsze spowoduje pełne
odwinięcie stosu: szczegóły zakończenia programu mogą zależeć od
implementacji.

\subsection{Przechwytywanie przez referencję i kolejność klauzul}

Przyjmujemy zwykle \texttt{const E\&}: pozwala to uniknąć niepotrzebnego
kopiowania i \emph{object slicing} przy przekazywaniu obiektu typu
pochodnego jako typu bazowego przez wartość.

\begin{lstlisting}
#include <exception>
#include <iostream>
#include <stdexcept>

void handle() {
    try {
        throw std::out_of_range("bad index");
    } catch (const std::out_of_range& e) {
        std::cout << "out_of_range: " << e.what() << '\n';
    } catch (const std::exception& e) {
        std::cout << "other std::exception: " << e.what() << '\n';
    } catch (...) {
        std::cout << "unknown exception\n";
    }
}
\end{lstlisting}

Klauzule są sprawdzane w kolejności występowania: najpierw umieszcza
się typy bardziej szczegółowe, a na końcu bazowe. Wyrażenie
\texttt{catch (...)} przechwytuje wyjątek dowolnego typu, ale nie
udostępnia bezpośrednio jego obiektu.

\paragraph{Ponowne zgłoszenie.}
\texttt{throw;} wewnątrz aktywnej obsługi ponownie zgłasza bieżący
wyjątek z zachowaniem jego dynamicznego typu. Natomiast
\texttt{throw e;} zgłasza nowe wyrażenie; w zależności od statycznego
typu \texttt{e} może prowadzić do kopiowania i obcięcia do klasy
bazowej.

\begin{lstlisting}
void log_and_rethrow() {
    try {
        throw std::runtime_error("disk error");
    } catch (const std::exception& e) {
        std::cerr << e.what() << '\n';
        throw;                  // keeps original exception object
        // throw e;              // may slice to std::exception
    }
}
\end{lstlisting}

\section{Odwijanie stosu i RAII}
\label{sec:cpp11-unwinding}

Najważniejsza reguła praktyczna brzmi: \emph{zwalnianie zasobu powinno
wynikać z destruktora obiektu zarządzającego zasobem}. Stanowi to
rozwinięcie idei RAII poznanej w rozdziale~\ref{chap:wskazniki-cpp}.

\begin{lstlisting}
#include <iostream>
#include <stdexcept>
#include <string>
#include <utility>

struct Trace {
    std::string name;
    explicit Trace(std::string n) : name(std::move(n)) {
        std::cout << "create " << name << '\n';
    }
    ~Trace() noexcept {
        std::cout << "destroy " << name << '\n';
    }
};

void third() {
    Trace c{"C"};
    throw std::runtime_error("failure");
}
void second() {
    Trace b{"B"};
    third();
}
void first() {
    Trace a{"A"};
    second();
}
// try { first(); } catch (const std::exception&) {}
// Destruction order: C, B, A.
\end{lstlisting}

Wyjątek powoduje zakończenie odpowiednich zakresów, a obiekty
z lokalnym automatycznym czasem przechowywania są niszczone w odwrotnej
kolejności ich skonstruowania. Jeżeli konstruktor zgłosi wyjątek,
destruktor \emph{całego, nieutworzonego obiektu} nie jest wywoływany,
ale niszczone są już w pełni skonstruowane podobiekty bazowe i składowe.

\subsection{Zarządzanie plikiem przy pomocy RAII}

\begin{lstlisting}
#include <cstdio>
#include <memory>
#include <stdexcept>

struct FileCloser {
    void operator()(std::FILE* file) const noexcept {
        if (file) std::fclose(file);
    }
};
using File = std::unique_ptr<std::FILE, FileCloser>;

File open_file(const char* name) {
    File file{std::fopen(name, "rb")};
    if (!file) throw std::runtime_error("cannot open file");
    return file;
}
\end{lstlisting}

Nawet gdy po otwarciu pliku dalsze operacje zgłoszą wyjątek,
\texttt{unique\_ptr} wykona deleter. Zauważmy, że
\texttt{std::fclose} może sygnalizować błąd przy zamykaniu; prosty
noexcept-deleter ignoruje ten kod błędu. Jeżeli informacja o błędzie
zamknięcia jest istotna, należy zapewnić osobną, jawną operację
finalizacji, której wynik można obsłużyć \emph{przed} destrukcją.

\section{Gwarancje bezpieczeństwa wyjątkowego}

\begin{center}\small
\begin{tabularx}{\textwidth}{@{}p{0.32\textwidth}X@{}}
\toprule
Gwarancja lub narzędzie & Skutek w razie wyjątku \\
\midrule
Brak wycieku & Zasoby pozostają poprawnie zwalniane, zwykle przez RAII. \\
Gwarancja podstawowa & Obiekty pozostają poprawne, choć ich stan może się zmienić. \\
Gwarancja silna & Operacja kończy się sukcesem albo pozostawia obserwowalny stan bez zmian. \\
Gwarancja braku wyjątków & Operacja nie pozwala wyjątkowi opuścić funkcji. \\
\code{noexcept} & Deklaruje gwarancję braku wyjątku; wyjątek opuszczający funkcję kończy program przez \code{std::terminate}. \\
\bottomrule
\end{tabularx}
\end{center}
\label{sec:cpp11-guarantees}

Bezpieczeństwo wyjątkowe nie oznacza tylko braku wycieków pamięci.
Załóżmy, że lista zawiera \texttt{[1, 2]} i dopisanie \texttt{3}
zgłosi wyjątek. Gwarancja podstawowa oznacza, że lista nadal nadaje
się do użycia, choć może mieć zmienioną zawartość. Gwarancja silna
oznacza, że nadal zawiera dokładnie \texttt{[1, 2]}. Gwarancja braku
wyjątków mówi tylko, że wyjątek nie opuści funkcji; jeśli operacja
ma dodatkowo zawsze osiągnąć cel, mówimy o gwarancji \emph{no-fail}.

Silna gwarancja nie oznacza automatycznie cofnięcia efektów
zewnętrznych, np. wysłanego pakietu czy wpisu w bazie danych.
Zakres transakcyjności trzeba sprecyzować w kontrakcie.

\subsection{Silna gwarancja dzięki operacji „przygotuj, a potem zatwierdź”}

\begin{lstlisting}
#include <utility>
#include <vector>

class Table {
    std::vector<int> data_;
public:
    explicit Table(std::vector<int> values)
        : data_(std::move(values)) {}

    void append_many(const std::vector<int>& more) {
        auto candidate = data_;          // may throw
        candidate.insert(candidate.end(),
                         more.begin(), more.end()); // may throw
        data_.swap(candidate);           // noexcept for std::vector<int>
    }

    const std::vector<int>& data() const noexcept { return data_; }
};
\end{lstlisting}

Jeżeli kopiowanie lub \texttt{insert} zawiedzie, \texttt{data\_}
pozostaje nietknięte. Dopiero bezwyjątkowe \texttt{swap} zatwierdza
przygotowaną wersję. Kosztem jest dodatkowa pamięć i kopiowanie;
silna gwarancja bywa więc świadomym wyborem projektowym, a nie darmową
optymalizacją.

\paragraph{Uwaga na aliasing.}
W przykładzie argument \texttt{more} może nawet odnosić się do
\texttt{data\_} obiektu \texttt{Table}: kopia \texttt{candidate}
powstaje przed modyfikacją oryginału. W kodzie operującym bezpośrednio
na tym samym kontenerze trzeba dodatkowo analizować ważność iteratorów
i referencji.

\subsection{Konstruktor i niezmienniki klasy}

Dobrze napisany konstruktor tworzy w pełni poprawny obiekt albo
zgłasza wyjątek. Bezpiecznie zarządzane składowe automatycznie
posprzątają po częściowo wykonanej inicjalizacji.

\begin{lstlisting}
#include <memory>
#include <stdexcept>

class Session {
    std::unique_ptr<int[]> buffer_;
    std::size_t size_;
public:
    explicit Session(std::size_t n)
        : buffer_(std::make_unique<int[]>(n)), size_(n) {
        if (n == 0) throw std::invalid_argument("empty session");
    }
    std::size_t size() const noexcept { return size_; }
};
\end{lstlisting}

Nie należy zastępować RAII ręcznie pisanymi \texttt{catch}, w których
sprząta się kolejne zasoby. Takie procedury są podatne na błędy,
zwłaszcza gdy konstruktor zdobywa wiele zasobów po kolei.

\section{\texttt{noexcept}: deklarowanie i sprawdzanie}
\label{sec:cpp11-noexcept}

\texttt{noexcept} w specyfikacji funkcji oznacza, że wyjątek nie może
opuścić jej ciała. Naruszenie tego wymagania kończy się wywołaniem
\texttt{std::terminate} (nie jest to „ciche zignorowanie” błędu).

\begin{lstlisting}
void close_handle() noexcept {
    // cleanup that must not let exceptions escape
}

void may_fail();
static_assert(noexcept(close_handle()));
static_assert(!noexcept(may_fail()));
\end{lstlisting}

W drugim użyciu \texttt{noexcept(expr)} jest \emph{operatorem}
zwracającym wartość logiczną w czasie kompilacji. Nie uruchamia on
samego wyrażenia.

\subsection{Warunkowe \texttt{noexcept}}

\begin{lstlisting}
#include <type_traits>
#include <utility>

template<class T>
void swap_values(T& a, T& b)
    noexcept(std::is_nothrow_move_constructible_v<T> &&
             std::is_nothrow_move_assignable_v<T>) {
    T tmp(std::move(a));
    a = std::move(b);
    b = std::move(tmp);
}
\end{lstlisting}

Jest to przykład \emph{warunkowej specyfikacji}. Nie oznacza,
że przedstawiony algorytm daje silną gwarancję, jeżeli przenoszenie
rzuca wyjątki; dla takich typów stan \texttt{a} lub \texttt{b} może
się zmienić przed niepowodzeniem. W kodzie generycznym trzeba
analizować nie tylko składnię, ale też gwarancje operacji typu \texttt{T}.

\subsection{Dlaczego \texttt{noexcept} jest ważne dla \texttt{vector}?}

Przy powiększaniu pojemności kontener może potrzebować przenieść
lub skopiować stare elementy do nowej pamięci. Dla typów kopiowalnych
konstrukcja kopiująca bywa bezpieczniejsza, jeśli konstruktor
przenoszący może zgłosić wyjątek.

\begin{lstlisting}
#include <type_traits>
#include <utility>

struct Item {
    Item() = default;
    Item(const Item&) {}              // copying available
    Item(Item&&) noexcept(false) {}   // move may throw
};

static_assert(std::is_copy_constructible_v<Item>);
static_assert(!std::is_nothrow_move_constructible_v<Item>);

// std::move_if_noexcept(x) usually gives const Item& here.
\end{lstlisting}

\texttt{std::move\_if\_noexcept} zwraca \texttt{T\&\&}, jeżeli
przenoszenie jest zadeklarowane jako bezwyjątkowe albo typ nie jest
kopiowalny; w przeciwnym razie zwraca \texttt{const T\&}.
Zasada ta pomaga kontenerom utrzymywać silną gwarancję w wielu,
ale nie we wszystkich przypadkach. Szczegóły są zależne od
konkretnej operacji i własności typu elementów.

\subsection{Destruktory i \texttt{std::terminate}}

Domyślnie destruktor jest niewyrzucający, o ile destruktory jego
składowych i klas bazowych nie wymagają innej specyfikacji.
Szczególnie podczas odwijania stosu wypuszczenie kolejnego wyjątku
z destruktora może spowodować \texttt{std::terminate}.

\begin{lstlisting}
struct SafeCleanup {
    ~SafeCleanup() noexcept {
        try {
            risky_cleanup();
        } catch (...) {
            // Log, record an error, or use another no-throw fallback.
            // Never let it escape this destructor.
        }
    }
    static void risky_cleanup();
};
\end{lstlisting}

Nie oznacza to, że każdą awarię należy ukryć. Jeśli operacja
finalizacji może się nie udać i użytkownik musi znać wynik, projektujmy
jawną metodę \texttt{close()} lub \texttt{commit()}, a destruktor
niech pełni rolę bezpiecznej ścieżki awaryjnej.

\section{Własne klasy wyjątków i hierarchia standardowa}
\label{sec:cpp11-types}

W nagłówkach \texttt{<stdexcept>}, \texttt{<exception>},
\texttt{<new>} i \texttt{<system\_error>} znajdziemy wiele gotowych
typów. W typowych zastosowaniach najczęściej spotkamy
\texttt{std::invalid\_argument}, \texttt{std::out\_of\_range},
\texttt{std::runtime\_error}, \texttt{std::bad\_alloc} oraz
\texttt{std::system\_error}.

\begin{lstlisting}
#include <cstddef>
#include <stdexcept>
#include <string>
#include <utility>

class ParseError : public std::runtime_error {
    std::size_t position_;
public:
    ParseError(std::size_t pos, std::string message)
        : std::runtime_error(std::move(message)),
          position_(pos) {}

    std::size_t position() const noexcept { return position_; }
};

int parse_digit(char c, std::size_t pos) {
    if (c < '0' || c > '9')
        throw ParseError(pos, "expected decimal digit");
    return c - '0';
}
\end{lstlisting}

Własna klasa może przechowywać dodatkowe informacje, np. identyfikator
operacji, pozycję w tekście lub kod błędu. Nie należy umieszczać tam
niepotrzebnie ogromnych danych. Metoda \texttt{what()} klasy
\texttt{std::exception} jest wirtualna i \texttt{noexcept}; jeśli
ją przesłaniamy, musimy zachować zgodny kontrakt i poprawny czas życia
zwracanego łańcucha.

\paragraph{Błędy systemowe.}
\texttt{std::system\_error} umożliwia przenoszenie kodu błędu
\texttt{std::error\_code} wraz z kategorią, np. błędu systemowego,
bez uzależniania całego API od przypadkowego tekstu \texttt{what()}.

\section{Wyjątki w funkcjach szablonowych i neutralność wyjątkowa}

Komponent generyczny powinien zwykle zachować niezmienniki i pozwolić
wyjątkom operacji wywołanych na typie \texttt{T} przejść do klienta,
zamiast pochłaniać je przez \texttt{catch (...)}.

\begin{lstlisting}
#include <algorithm>
#include <vector>

template<class T, class Compare>
void sort_elements(std::vector<T>& xs, Compare comp) {
    std::sort(xs.begin(), xs.end(), comp);
    // An exception from comp may propagate to the caller.
}
\end{lstlisting}

Należy jednak odróżnić \emph{neutralność wyjątkową} od silnej
gwarancji. Jeśli komparator zgłosi wyjątek w trakcie sortowania,
kontener pozostaje poprawnym obiektem, ale jego kolejność nie musi
odpowiadać stanowi sprzed wywołania.

\section{Przenoszenie wyjątków między funkcjami i wątkami}
\label{sec:cpp11-exception-ptr}

Ta i dwie kolejne sekcje pokazują zastosowania bardziej zaawansowane.
Do podstaw wystarczą \texttt{try}/\texttt{catch}, RAII oraz gwarancje
stanu po niepowodzeniu.

Typ \texttt{std::exception\_ptr} pozwala zachować wyjątek poza blokiem
\texttt{catch}, a później zgłosić go ponownie. Jest przydatny
szczególnie w programowaniu współbieżnym.

\begin{lstlisting}
#include <exception>
#include <iostream>
#include <stdexcept>
#include <string>
#include <thread>

std::exception_ptr error;
std::thread worker([&] {
    try {
        throw std::runtime_error("worker failed");
    } catch (...) {
        error = std::current_exception();
    }
});
worker.join();

if (error) {
    std::rethrow_exception(error);
}
\end{lstlisting}

Zmienna \texttt{error} nie jest tutaj współdzielona jednocześnie przez
dwa aktywne dostępy: zapis zachodzi w wątku roboczym, a odczyt
dopiero po \texttt{join()}, które zapewnia wymaganą synchronizację.
Nie wolno dopuścić, aby wyjątek opuścił funkcję startową zwykłego
\texttt{std::thread}: typowym skutkiem jest \texttt{std::terminate}.

\subsection{\texttt{std::future} przechowuje wyjątek}

Wspomniane w rozdziale o współbieżności \texttt{std::async} i
\texttt{std::future} ułatwiają propagowanie niepowodzenia do
odbiorcy wyniku.

\begin{lstlisting}
#include <future>
#include <stdexcept>

auto fut = std::async(std::launch::async, []() -> int {
    throw std::runtime_error("calculation failed");
});
try {
    (void)fut.get();  // rethrows the stored exception
} catch (const std::runtime_error&) {
    // report the failure in the receiving thread
}
\end{lstlisting}

Przez \texttt{get()} wynik (lub wyjątek) trafia do wątku odbiorcy.
W \texttt{std::future} wynik można pobrać tylko raz; inne własności
tego mechanizmu omówiliśmy w rozdziale o asynchroniczności.

\section{Wyjątki zagnieżdżone i dodatkowy kontekst}
\label{sec:cpp11-nested}

Gdy funkcja wyższego poziomu rozumie więcej niż warstwa niższa, może
wzbogacić informację o błędzie i zachować oryginalną przyczynę.

\begin{lstlisting}
#include <exception>
#include <stdexcept>

void load_config() {
    try {
        throw std::runtime_error("invalid token");
    } catch (...) {
        std::throw_with_nested(
            std::runtime_error("cannot load configuration"));
    }
}

void print_nested(const std::exception& e, int depth = 0) {
    std::cerr << std::string(depth * 2, ' ') << e.what() << '\n';
    try {
        std::rethrow_if_nested(e);
    } catch (const std::exception& inner) {
        print_nested(inner, depth + 1);
    } catch (...) {
        std::cerr << "nested non-standard exception\n";
    }
}
\end{lstlisting}

Takie opakowywanie wyjątków nie jest tym
samym co zgłoszenie \texttt{throw e;}: nowy wyjątek przechowuje
kontekst, a przyczynę można odzyskać przez \texttt{rethrow\_if\_nested}.

\section{\texttt{std::uncaught\_exceptions} i strażnicy zakresu}

Funkcja \texttt{std::uncaught\_exceptions()} (C++17) zwraca liczbę
aktualnie nieobsługiwanych obiektów wyjątków w bieżącym wątku.
Można jej użyć do wykrycia, czy wyjście z zakresu następuje wskutek
odwijania stosu.

\begin{lstlisting}
#include <exception>

class FailureGuard {
    int initial_ = std::uncaught_exceptions();
    int& failure_count_;
public:
    explicit FailureGuard(int& count) noexcept : failure_count_(count) {}
    FailureGuard(const FailureGuard&) = delete;
    FailureGuard& operator=(const FailureGuard&) = delete;

    ~FailureGuard() noexcept {
        if (std::uncaught_exceptions() > initial_) {
            ++failure_count_; // no-throw bookkeeping
        }
    }
};
\end{lstlisting}

Porównanie z \emph{początkową liczbą}, a nie tylko z zerem, jest
ważne przy zagnieżdżonych wyjątkach. Strażnik nie daje jednak
pełnego, transakcyjnego mechanizmu cofania zmian. W projektach
korzystających z C++20 można napisać własny scope guard, ale trzeba
uważnie określić, czy przechowuje funkcję przez kopię, przez
przeniesienie czy przez referencję i czy jej wywołanie może rzucać.

\section{Wyjątki a alternatywne modele obsługi błędów}
\label{sec:cpp11-alternatives}

\subsection{\texttt{std::optional}: brak wyniku nie musi być błędem}

\begin{lstlisting}
#include <optional>
#include <string_view>

std::optional<std::size_t> find_colon(std::string_view s) {
    const auto i = s.find(':');
    if (i == std::string_view::npos) return std::nullopt;
    return i;
}
\end{lstlisting}

Brak dwukropka może być normalnym scenariuszem. Zgłaszanie wyjątku
przy każdym nieudanym wyszukiwaniu bywa niepotrzebne.

\subsection{\texttt{std::expected<T,E>} — C++23, nie C++20}

\texttt{std::expected} reprezentuje \emph{albo wartość typu T, albo
błąd typu E}. Został ustandaryzowany w C++23 i nie należy zakładać,
że nagłówek \texttt{<expected>} jest dostępny w kompilacji C++20.

\begin{lstlisting}
// This example requires C++23 and <expected>.
#include <expected>
#include <string>
#include <string_view>

std::expected<int, std::string>
parse_positive(std::string_view text) {
    if (text == "1") return 1;
    if (text == "2") return 2;
    return std::unexpected("expected 1 or 2");
}

void use_expected() {
    auto result = parse_positive("3");
    if (!result) {
        const std::string& error = result.error();
        (void)error;
    }
}
\end{lstlisting}

Dzięki temu błędy oczekiwane przez API są jawne w sygnaturze.
\texttt{std::expected} nie eliminuje wszystkich wyjątków: konstrukcja
zawartego obiektu czy inne operacje nadal mogą się nie udać.
W C++20 podobny kontrakt można zrealizować własnym typem
\texttt{Result}, \texttt{std::variant<T,E>} lub odpowiednią biblioteką.

\subsection{Kody błędów, wyjątki i wydajność}

W konwencjonalnych implementacjach C++ obsługa sukcesu na ścieżce,
na której żaden wyjątek nie wystąpił, może być bardzo tania, ale samo
zgłoszenie i obsłużenie wyjątku zwykle kosztuje więcej niż zwrócenie
prostego kodu błędu. Nie należy jednak zamieniać każdego \texttt{throw}
na kod zwrotny tylko na podstawie intuicji: warto mierzyć rzeczywistą
częstotliwość błędów, profilować aplikację i uwzględniać prostotę
kontraktów.

\section{Typowe błędy i dobre praktyki}
\label{sec:cpp11-guidelines}

\begin{enumerate}
\item \textbf{Nie rzucaj wskaźnikami do lokalnych obiektów.}
      Preferuj wyjątki będące obiektami, zwykle pochodnymi od
      \texttt{std::exception}.
\item \textbf{Przechwytuj wyjątki polimorficzne przez
      \texttt{const\&}.} Unikaj obcinania obiektu (\emph{slicing}).
\item \textbf{Do ponownego zgłoszenia używaj \texttt{throw;}.}
      Nie zastępuj go bez namysłu \texttt{throw e;}.
\item \textbf{Stosuj RAII.} Surowe wywołania \texttt{new}/\texttt{delete}
      rozdzielone operacjami rzucającymi wyjątki to częsta przyczyna
      wycieków pamięci.
\item \textbf{Nie wypuszczaj wyjątków z destruktorów.}
      Błędy finalizacji zgłaszaj jawnie w odpowiednio zaprojektowanym API.
\item \textbf{Nie oznaczaj funkcji \texttt{noexcept} bez analizy.}
      To obietnica z konsekwencją w postaci \texttt{std::terminate}.
\item \textbf{Nie ukrywaj błędów w pustym \texttt{catch (...)}.}
      Obsłuż wyjątek, zapisz informację diagnostyczną albo przekaż go wyżej.
\item \textbf{Określ gwarancję niepowodzenia.} Dla operacji mutujących
      dane rozważ silną gwarancję i strategię „przygotuj–zatwierdź”.
\item \textbf{Złap wyjątek w funkcji wątku.}
      Nie pozwól mu opuścić funkcji startowej \texttt{std::thread}.
      Do przekazania błędu użyj \texttt{std::exception\_ptr},
      \texttt{promise/future} lub API wyższego poziomu.
\item \textbf{Wybierz właściwy model błędu.} Brak zwykłego wyniku,
      przewidywalna porażka operacji i naruszenie kontraktu to różne
      sytuacje.
\end{enumerate}

\section{Przykład końcowy: transakcyjna aktualizacja danych}

Poniższy program ilustruje współpracę wyjątku, RAII i silnej gwarancji.
Załóżmy, że aktualizacja dodaje wartości do rejestru, ale w razie błędu
rejestr ma pozostać bez zmian.

\begin{lstlisting}
#include <iostream>
#include <stdexcept>
#include <utility>
#include <vector>

class Register {
    std::vector<int> values_;
public:
    Register() : values_{1, 2, 3} {}

    void add_checked(const std::vector<int>& numbers) {
        std::vector<int> proposed = values_;
        for (int x : numbers) {
            if (x < 0) {
                throw std::invalid_argument("negative input");
            }
            proposed.push_back(x);
        }
        values_.swap(proposed);
    }

    const std::vector<int>& values() const noexcept {
        return values_;
    }
};

int main() {
    Register r;
    try {
        r.add_checked({4, -5, 6});
    } catch (const std::invalid_argument& e) {
        std::cerr << e.what() << '\n';
    }
    // Still {1, 2, 3}; no partial update.
    for (int x : r.values()) std::cout << x << ' ';
}
\end{lstlisting}

Najpierw powstaje kopia danych. Sprawdzanie i alokowanie nowej pamięci
mogą zgłaszać wyjątki, lecz działają na kopii. Po pomyślnym wykonaniu
następuje zatwierdzenie przez \texttt{swap}. Jest to wzorzec, który
można stosować również w rozbudowanych klasach, przy świadomym
uwzględnieniu kosztów i własności ich składowych.

\section{Podsumowanie rozdziału}

W praktyce mechanizm wyjątków jest najbardziej użyteczny wtedy,
gdy łączy się go z dobrym modelem własności zasobów i z precyzyjnie
określonymi gwarancjami operacji. \texttt{throw}, \texttt{catch}
i \texttt{noexcept} to tylko składnia; o poprawności projektu
decydują RAII, niezmienniki obiektów, semantyka transakcyjna oraz
świadome oddzielanie błędów wyjątkowych od zwykłych wariantów wyniku.



\chapter{Wyjątki i bezpieczeństwo wyjątkowe w C++}
\label{chap:cpp-exceptions}

\section{Po co C++ ma wyjątki?}

Mechanizm wyjątków oddziela kod realizujący zwykły przebieg algorytmu
od kodu reagującego na niepowodzenia. Jeżeli funkcja nie może spełnić
swojego kontraktu, może zgłosić wyjątek poleceniem \texttt{throw}.
Wówczas wykonanie jest przekazywane do odpowiedniego bloku
\texttt{catch}; po drodze niszczone są obiekty automatyczne, których
czas życia dobiegł końca. To ostatnie zjawisko nazywamy
\emph{odwijaniem stosu} (\emph{stack unwinding}).

Najprostszy schemat wygląda tak: \texttt{throw} zgłasza problem,
a pasujący \texttt{catch} go odbiera. Kod po \texttt{throw} w tym samym
bloku nie zostanie wykonany:
\begin{lstlisting}
try {
    throw std::runtime_error("bad input");
} catch (const std::runtime_error& error) {
    std::cerr << error.what() << '\n';
}
\end{lstlisting}
W samodzielnym programie ten fragment wymaga nagłówków
\texttt{<iostream>} i \texttt{<stdexcept>}.

Z perspektywy projektowania API warto odróżniać:
\begin{itemize}
\item \emph{zwykły wynik} (w tym informację, że elementu nie znaleziono),
      dla którego często lepszy jest \texttt{std::optional};
\item \emph{błąd operacji, który wywołujący powinien świadomie obsłużyć},
      reprezentowany np. przez kod błędu lub, od C++23, \texttt{std::expected};
\item \emph{niepowodzenie uniemożliwiające dotrzymanie kontraktu},
      dla którego wyjątek bywa naturalnym sposobem propagacji.
\end{itemize}
Nie ma uniwersalnej reguły nakazującej używanie wyjątków wszędzie.
Znaczenie mają kontrakty funkcji, konwencje projektu, ograniczenia
środowiska oraz wymagania dotyczące opóźnień wykonania.

\section{Składnia \texttt{try}, \texttt{throw} i \texttt{catch}}

\begin{center}\small
\begin{tabularx}{\textwidth}{@{}p{0.32\textwidth}X@{}}
\toprule
Konstrukcja & Za co odpowiada \\
\midrule
\code{try} & Obejmuje blok, z którego wyjątki mogą zostać przechwycone przez następujące bloki \code{catch}. \\
\code{throw value} & Tworzy obiekt wyjątku i rozpoczyna poszukiwanie pasującej obsługi. \\
\texttt{catch(const E\& e)} & Przechwytuje wyjątek zgodnego typu bez kopiowania i bez obcinania części pochodnej. \\
\code{catch(...)} & Przechwytuje dowolny wyjątek, także nieznanego typu. \\
\code{throw;} & Wewnątrz obsługi ponownie zgłasza bieżący wyjątek, zachowując jego typ dynamiczny. \\
\code{std::exception::what()} & Zwraca tekstowy opis błędu; służy do diagnostyki, nie do sterowania logiką. \\
\bottomrule
\end{tabularx}
\end{center}
\label{sec:cpp11-basics}

\begin{lstlisting}
#include <iostream>
#include <stdexcept>

double divide(double a, double b) {
    if (b == 0.0)
        throw std::invalid_argument("zero denominator");
    return a / b;
}

int main() {
    try {
        std::cout << divide(10.0, 2.0) << '\n';
        std::cout << divide(10.0, 0.0) << '\n';
        std::cout << "not executed\n";
    } catch (const std::invalid_argument& e) {
        std::cerr << "invalid input: " << e.what() << '\n';
    }
}
\end{lstlisting}

Instrukcja \texttt{throw E\{...\}} tworzy obiekt wyjątku. Obsługa
szuka pasującej klauzuli \texttt{catch}; jeśli wyjątek nie zostanie
obsłużony, program kończy działanie poprzez \texttt{std::terminate}.
Nie należy zakładać, że nieobsłużony wyjątek zawsze spowoduje pełne
odwinięcie stosu: szczegóły zakończenia programu mogą zależeć od
implementacji.

\subsection{Przechwytywanie przez referencję i kolejność klauzul}

Przyjmujemy zwykle \texttt{const E\&}: pozwala to uniknąć niepotrzebnego
kopiowania i \emph{object slicing} przy przekazywaniu obiektu typu
pochodnego jako typu bazowego przez wartość.

\begin{lstlisting}
#include <exception>
#include <iostream>
#include <stdexcept>

void handle() {
    try {
        throw std::out_of_range("bad index");
    } catch (const std::out_of_range& e) {
        std::cout << "out_of_range: " << e.what() << '\n';
    } catch (const std::exception& e) {
        std::cout << "other std::exception: " << e.what() << '\n';
    } catch (...) {
        std::cout << "unknown exception\n";
    }
}
\end{lstlisting}

Klauzule są sprawdzane w kolejności występowania: najpierw umieszcza
się typy bardziej szczegółowe, a na końcu bazowe. Wyrażenie
\texttt{catch (...)} przechwytuje wyjątek dowolnego typu, ale nie
udostępnia bezpośrednio jego obiektu.

\paragraph{Ponowne zgłoszenie.}
\texttt{throw;} wewnątrz aktywnej obsługi ponownie zgłasza bieżący
wyjątek z zachowaniem jego dynamicznego typu. Natomiast
\texttt{throw e;} zgłasza nowe wyrażenie; w zależności od statycznego
typu \texttt{e} może prowadzić do kopiowania i obcięcia do klasy
bazowej.

\begin{lstlisting}
void log_and_rethrow() {
    try {
        throw std::runtime_error("disk error");
    } catch (const std::exception& e) {
        std::cerr << e.what() << '\n';
        throw;                  // keeps original exception object
        // throw e;              // may slice to std::exception
    }
}
\end{lstlisting}

\section{Odwijanie stosu i RAII}
\label{sec:cpp11-unwinding}

Najważniejsza reguła praktyczna brzmi: \emph{zwalnianie zasobu powinno
wynikać z destruktora obiektu zarządzającego zasobem}. Stanowi to
rozwinięcie idei RAII poznanej w rozdziale~\ref{chap:wskazniki-cpp}.

\begin{lstlisting}
#include <iostream>
#include <stdexcept>
#include <string>
#include <utility>

struct Trace {
    std::string name;
    explicit Trace(std::string n) : name(std::move(n)) {
        std::cout << "create " << name << '\n';
    }
    ~Trace() noexcept {
        std::cout << "destroy " << name << '\n';
    }
};

void third() {
    Trace c{"C"};
    throw std::runtime_error("failure");
}
void second() {
    Trace b{"B"};
    third();
}
void first() {
    Trace a{"A"};
    second();
}
// try { first(); } catch (const std::exception&) {}
// Destruction order: C, B, A.
\end{lstlisting}

Wyjątek powoduje zakończenie odpowiednich zakresów, a obiekty
z lokalnym automatycznym czasem przechowywania są niszczone w odwrotnej
kolejności ich skonstruowania. Jeżeli konstruktor zgłosi wyjątek,
destruktor \emph{całego, nieutworzonego obiektu} nie jest wywoływany,
ale niszczone są już w pełni skonstruowane podobiekty bazowe i składowe.

\subsection{Zarządzanie plikiem przy pomocy RAII}

\begin{lstlisting}
#include <cstdio>
#include <memory>
#include <stdexcept>

struct FileCloser {
    void operator()(std::FILE* file) const noexcept {
        if (file) std::fclose(file);
    }
};
using File = std::unique_ptr<std::FILE, FileCloser>;

File open_file(const char* name) {
    File file{std::fopen(name, "rb")};
    if (!file) throw std::runtime_error("cannot open file");
    return file;
}
\end{lstlisting}

Nawet gdy po otwarciu pliku dalsze operacje zgłoszą wyjątek,
\texttt{unique\_ptr} wykona deleter. Zauważmy, że
\texttt{std::fclose} może sygnalizować błąd przy zamykaniu; prosty
noexcept-deleter ignoruje ten kod błędu. Jeżeli informacja o błędzie
zamknięcia jest istotna, należy zapewnić osobną, jawną operację
finalizacji, której wynik można obsłużyć \emph{przed} destrukcją.

\section{Gwarancje bezpieczeństwa wyjątkowego}

\begin{center}\small
\begin{tabularx}{\textwidth}{@{}p{0.32\textwidth}X@{}}
\toprule
Gwarancja lub narzędzie & Skutek w razie wyjątku \\
\midrule
Brak wycieku & Zasoby pozostają poprawnie zwalniane, zwykle przez RAII. \\
Gwarancja podstawowa & Obiekty pozostają poprawne, choć ich stan może się zmienić. \\
Gwarancja silna & Operacja kończy się sukcesem albo pozostawia obserwowalny stan bez zmian. \\
Gwarancja braku wyjątków & Operacja nie pozwala wyjątkowi opuścić funkcji. \\
\code{noexcept} & Deklaruje gwarancję braku wyjątku; wyjątek opuszczający funkcję kończy program przez \code{std::terminate}. \\
\bottomrule
\end{tabularx}
\end{center}
\label{sec:cpp11-guarantees}

Bezpieczeństwo wyjątkowe nie oznacza tylko braku wycieków pamięci.
Załóżmy, że lista zawiera \texttt{[1, 2]} i dopisanie \texttt{3}
zgłosi wyjątek. Gwarancja podstawowa oznacza, że lista nadal nadaje
się do użycia, choć może mieć zmienioną zawartość. Gwarancja silna
oznacza, że nadal zawiera dokładnie \texttt{[1, 2]}. Gwarancja braku
wyjątków mówi tylko, że wyjątek nie opuści funkcji; jeśli operacja
ma dodatkowo zawsze osiągnąć cel, mówimy o gwarancji \emph{no-fail}.

Silna gwarancja nie oznacza automatycznie cofnięcia efektów
zewnętrznych, np. wysłanego pakietu czy wpisu w bazie danych.
Zakres transakcyjności trzeba sprecyzować w kontrakcie.

\subsection{Silna gwarancja dzięki operacji „przygotuj, a potem zatwierdź”}

\begin{lstlisting}
#include <utility>
#include <vector>

class Table {
    std::vector<int> data_;
public:
    explicit Table(std::vector<int> values)
        : data_(std::move(values)) {}

    void append_many(const std::vector<int>& more) {
        auto candidate = data_;          // may throw
        candidate.insert(candidate.end(),
                         more.begin(), more.end()); // may throw
        data_.swap(candidate);           // noexcept for std::vector<int>
    }

    const std::vector<int>& data() const noexcept { return data_; }
};
\end{lstlisting}

Jeżeli kopiowanie lub \texttt{insert} zawiedzie, \texttt{data\_}
pozostaje nietknięte. Dopiero bezwyjątkowe \texttt{swap} zatwierdza
przygotowaną wersję. Kosztem jest dodatkowa pamięć i kopiowanie;
silna gwarancja bywa więc świadomym wyborem projektowym, a nie darmową
optymalizacją.

\paragraph{Uwaga na aliasing.}
W przykładzie argument \texttt{more} może nawet odnosić się do
\texttt{data\_} obiektu \texttt{Table}: kopia \texttt{candidate}
powstaje przed modyfikacją oryginału. W kodzie operującym bezpośrednio
na tym samym kontenerze trzeba dodatkowo analizować ważność iteratorów
i referencji.

\subsection{Konstruktor i niezmienniki klasy}

Dobrze napisany konstruktor tworzy w pełni poprawny obiekt albo
zgłasza wyjątek. Bezpiecznie zarządzane składowe automatycznie
posprzątają po częściowo wykonanej inicjalizacji.

\begin{lstlisting}
#include <memory>
#include <stdexcept>

class Session {
    std::unique_ptr<int[]> buffer_;
    std::size_t size_;
public:
    explicit Session(std::size_t n)
        : buffer_(std::make_unique<int[]>(n)), size_(n) {
        if (n == 0) throw std::invalid_argument("empty session");
    }
    std::size_t size() const noexcept { return size_; }
};
\end{lstlisting}

Nie należy zastępować RAII ręcznie pisanymi \texttt{catch}, w których
sprząta się kolejne zasoby. Takie procedury są podatne na błędy,
zwłaszcza gdy konstruktor zdobywa wiele zasobów po kolei.

\section{\texttt{noexcept}: deklarowanie i sprawdzanie}
\label{sec:cpp11-noexcept}

\texttt{noexcept} w specyfikacji funkcji oznacza, że wyjątek nie może
opuścić jej ciała. Naruszenie tego wymagania kończy się wywołaniem
\texttt{std::terminate} (nie jest to „ciche zignorowanie” błędu).

\begin{lstlisting}
void close_handle() noexcept {
    // cleanup that must not let exceptions escape
}

void may_fail();
static_assert(noexcept(close_handle()));
static_assert(!noexcept(may_fail()));
\end{lstlisting}

W drugim użyciu \texttt{noexcept(expr)} jest \emph{operatorem}
zwracającym wartość logiczną w czasie kompilacji. Nie uruchamia on
samego wyrażenia.

\subsection{Warunkowe \texttt{noexcept}}

\begin{lstlisting}
#include <type_traits>
#include <utility>

template<class T>
void swap_values(T& a, T& b)
    noexcept(std::is_nothrow_move_constructible_v<T> &&
             std::is_nothrow_move_assignable_v<T>) {
    T tmp(std::move(a));
    a = std::move(b);
    b = std::move(tmp);
}
\end{lstlisting}

Jest to przykład \emph{warunkowej specyfikacji}. Nie oznacza,
że przedstawiony algorytm daje silną gwarancję, jeżeli przenoszenie
rzuca wyjątki; dla takich typów stan \texttt{a} lub \texttt{b} może
się zmienić przed niepowodzeniem. W kodzie generycznym trzeba
analizować nie tylko składnię, ale też gwarancje operacji typu \texttt{T}.

\subsection{Dlaczego \texttt{noexcept} jest ważne dla \texttt{vector}?}

Przy powiększaniu pojemności kontener może potrzebować przenieść
lub skopiować stare elementy do nowej pamięci. Dla typów kopiowalnych
konstrukcja kopiująca bywa bezpieczniejsza, jeśli konstruktor
przenoszący może zgłosić wyjątek.

\begin{lstlisting}
#include <type_traits>
#include <utility>

struct Item {
    Item() = default;
    Item(const Item&) {}              // copying available
    Item(Item&&) noexcept(false) {}   // move may throw
};

static_assert(std::is_copy_constructible_v<Item>);
static_assert(!std::is_nothrow_move_constructible_v<Item>);

// std::move_if_noexcept(x) usually gives const Item& here.
\end{lstlisting}

\texttt{std::move\_if\_noexcept} zwraca \texttt{T\&\&}, jeżeli
przenoszenie jest zadeklarowane jako bezwyjątkowe albo typ nie jest
kopiowalny; w przeciwnym razie zwraca \texttt{const T\&}.
Zasada ta pomaga kontenerom utrzymywać silną gwarancję w wielu,
ale nie we wszystkich przypadkach. Szczegóły są zależne od
konkretnej operacji i własności typu elementów.

\subsection{Destruktory i \texttt{std::terminate}}

Domyślnie destruktor jest niewyrzucający, o ile destruktory jego
składowych i klas bazowych nie wymagają innej specyfikacji.
Szczególnie podczas odwijania stosu wypuszczenie kolejnego wyjątku
z destruktora może spowodować \texttt{std::terminate}.

\begin{lstlisting}
struct SafeCleanup {
    ~SafeCleanup() noexcept {
        try {
            risky_cleanup();
        } catch (...) {
            // Log, record an error, or use another no-throw fallback.
            // Never let it escape this destructor.
        }
    }
    static void risky_cleanup();
};
\end{lstlisting}

Nie oznacza to, że każdą awarię należy ukryć. Jeśli operacja
finalizacji może się nie udać i użytkownik musi znać wynik, projektujmy
jawną metodę \texttt{close()} lub \texttt{commit()}, a destruktor
niech pełni rolę bezpiecznej ścieżki awaryjnej.

\section{Własne klasy wyjątków i hierarchia standardowa}
\label{sec:cpp11-types}

W nagłówkach \texttt{<stdexcept>}, \texttt{<exception>},
\texttt{<new>} i \texttt{<system\_error>} znajdziemy wiele gotowych
typów. W typowych zastosowaniach najczęściej spotkamy
\texttt{std::invalid\_argument}, \texttt{std::out\_of\_range},
\texttt{std::runtime\_error}, \texttt{std::bad\_alloc} oraz
\texttt{std::system\_error}.

\begin{lstlisting}
#include <cstddef>
#include <stdexcept>
#include <string>
#include <utility>

class ParseError : public std::runtime_error {
    std::size_t position_;
public:
    ParseError(std::size_t pos, std::string message)
        : std::runtime_error(std::move(message)),
          position_(pos) {}

    std::size_t position() const noexcept { return position_; }
};

int parse_digit(char c, std::size_t pos) {
    if (c < '0' || c > '9')
        throw ParseError(pos, "expected decimal digit");
    return c - '0';
}
\end{lstlisting}

Własna klasa może przechowywać dodatkowe informacje, np. identyfikator
operacji, pozycję w tekście lub kod błędu. Nie należy umieszczać tam
niepotrzebnie ogromnych danych. Metoda \texttt{what()} klasy
\texttt{std::exception} jest wirtualna i \texttt{noexcept}; jeśli
ją przesłaniamy, musimy zachować zgodny kontrakt i poprawny czas życia
zwracanego łańcucha.

\paragraph{Błędy systemowe.}
\texttt{std::system\_error} umożliwia przenoszenie kodu błędu
\texttt{std::error\_code} wraz z kategorią, np. błędu systemowego,
bez uzależniania całego API od przypadkowego tekstu \texttt{what()}.

\section{Wyjątki w funkcjach szablonowych i neutralność wyjątkowa}

Komponent generyczny powinien zwykle zachować niezmienniki i pozwolić
wyjątkom operacji wywołanych na typie \texttt{T} przejść do klienta,
zamiast pochłaniać je przez \texttt{catch (...)}.

\begin{lstlisting}
#include <algorithm>
#include <vector>

template<class T, class Compare>
void sort_elements(std::vector<T>& xs, Compare comp) {
    std::sort(xs.begin(), xs.end(), comp);
    // An exception from comp may propagate to the caller.
}
\end{lstlisting}

Należy jednak odróżnić \emph{neutralność wyjątkową} od silnej
gwarancji. Jeśli komparator zgłosi wyjątek w trakcie sortowania,
kontener pozostaje poprawnym obiektem, ale jego kolejność nie musi
odpowiadać stanowi sprzed wywołania.

\section{Przenoszenie wyjątków między funkcjami i wątkami}
\label{sec:cpp11-exception-ptr}

Ta i dwie kolejne sekcje pokazują zastosowania bardziej zaawansowane.
Do podstaw wystarczą \texttt{try}/\texttt{catch}, RAII oraz gwarancje
stanu po niepowodzeniu.

Typ \texttt{std::exception\_ptr} pozwala zachować wyjątek poza blokiem
\texttt{catch}, a później zgłosić go ponownie. Jest przydatny
szczególnie w programowaniu współbieżnym.

\begin{lstlisting}
#include <exception>
#include <iostream>
#include <stdexcept>
#include <string>
#include <thread>

std::exception_ptr error;
std::thread worker([&] {
    try {
        throw std::runtime_error("worker failed");
    } catch (...) {
        error = std::current_exception();
    }
});
worker.join();

if (error) {
    std::rethrow_exception(error);
}
\end{lstlisting}

Zmienna \texttt{error} nie jest tutaj współdzielona jednocześnie przez
dwa aktywne dostępy: zapis zachodzi w wątku roboczym, a odczyt
dopiero po \texttt{join()}, które zapewnia wymaganą synchronizację.
Nie wolno dopuścić, aby wyjątek opuścił funkcję startową zwykłego
\texttt{std::thread}: typowym skutkiem jest \texttt{std::terminate}.

\subsection{\texttt{std::future} przechowuje wyjątek}

Wspomniane w rozdziale o współbieżności \texttt{std::async} i
\texttt{std::future} ułatwiają propagowanie niepowodzenia do
odbiorcy wyniku.

\begin{lstlisting}
#include <future>
#include <stdexcept>

auto fut = std::async(std::launch::async, []() -> int {
    throw std::runtime_error("calculation failed");
});
try {
    (void)fut.get();  // rethrows the stored exception
} catch (const std::runtime_error&) {
    // report the failure in the receiving thread
}
\end{lstlisting}

Przez \texttt{get()} wynik (lub wyjątek) trafia do wątku odbiorcy.
W \texttt{std::future} wynik można pobrać tylko raz; inne własności
tego mechanizmu omówiliśmy w rozdziale o asynchroniczności.

\section{Wyjątki zagnieżdżone i dodatkowy kontekst}
\label{sec:cpp11-nested}

Gdy funkcja wyższego poziomu rozumie więcej niż warstwa niższa, może
wzbogacić informację o błędzie i zachować oryginalną przyczynę.

\begin{lstlisting}
#include <exception>
#include <stdexcept>

void load_config() {
    try {
        throw std::runtime_error("invalid token");
    } catch (...) {
        std::throw_with_nested(
            std::runtime_error("cannot load configuration"));
    }
}

void print_nested(const std::exception& e, int depth = 0) {
    std::cerr << std::string(depth * 2, ' ') << e.what() << '\n';
    try {
        std::rethrow_if_nested(e);
    } catch (const std::exception& inner) {
        print_nested(inner, depth + 1);
    } catch (...) {
        std::cerr << "nested non-standard exception\n";
    }
}
\end{lstlisting}

Takie opakowywanie wyjątków nie jest tym
samym co zgłoszenie \texttt{throw e;}: nowy wyjątek przechowuje
kontekst, a przyczynę można odzyskać przez \texttt{rethrow\_if\_nested}.

\section{\texttt{std::uncaught\_exceptions} i strażnicy zakresu}

Funkcja \texttt{std::uncaught\_exceptions()} (C++17) zwraca liczbę
aktualnie nieobsługiwanych obiektów wyjątków w bieżącym wątku.
Można jej użyć do wykrycia, czy wyjście z zakresu następuje wskutek
odwijania stosu.

\begin{lstlisting}
#include <exception>

class FailureGuard {
    int initial_ = std::uncaught_exceptions();
    int& failure_count_;
public:
    explicit FailureGuard(int& count) noexcept : failure_count_(count) {}
    FailureGuard(const FailureGuard&) = delete;
    FailureGuard& operator=(const FailureGuard&) = delete;

    ~FailureGuard() noexcept {
        if (std::uncaught_exceptions() > initial_) {
            ++failure_count_; // no-throw bookkeeping
        }
    }
};
\end{lstlisting}

Porównanie z \emph{początkową liczbą}, a nie tylko z zerem, jest
ważne przy zagnieżdżonych wyjątkach. Strażnik nie daje jednak
pełnego, transakcyjnego mechanizmu cofania zmian. W projektach
korzystających z C++20 można napisać własny scope guard, ale trzeba
uważnie określić, czy przechowuje funkcję przez kopię, przez
przeniesienie czy przez referencję i czy jej wywołanie może rzucać.

\section{Wyjątki a alternatywne modele obsługi błędów}
\label{sec:cpp11-alternatives}

\subsection{\texttt{std::optional}: brak wyniku nie musi być błędem}

\begin{lstlisting}
#include <optional>
#include <string_view>

std::optional<std::size_t> find_colon(std::string_view s) {
    const auto i = s.find(':');
    if (i == std::string_view::npos) return std::nullopt;
    return i;
}
\end{lstlisting}

Brak dwukropka może być normalnym scenariuszem. Zgłaszanie wyjątku
przy każdym nieudanym wyszukiwaniu bywa niepotrzebne.

\subsection{\texttt{std::expected<T,E>} — C++23, nie C++20}

\texttt{std::expected} reprezentuje \emph{albo wartość typu T, albo
błąd typu E}. Został ustandaryzowany w C++23 i nie należy zakładać,
że nagłówek \texttt{<expected>} jest dostępny w kompilacji C++20.

\begin{lstlisting}
// This example requires C++23 and <expected>.
#include <expected>
#include <string>
#include <string_view>

std::expected<int, std::string>
parse_positive(std::string_view text) {
    if (text == "1") return 1;
    if (text == "2") return 2;
    return std::unexpected("expected 1 or 2");
}

void use_expected() {
    auto result = parse_positive("3");
    if (!result) {
        const std::string& error = result.error();
        (void)error;
    }
}
\end{lstlisting}

Dzięki temu błędy oczekiwane przez API są jawne w sygnaturze.
\texttt{std::expected} nie eliminuje wszystkich wyjątków: konstrukcja
zawartego obiektu czy inne operacje nadal mogą się nie udać.
W C++20 podobny kontrakt można zrealizować własnym typem
\texttt{Result}, \texttt{std::variant<T,E>} lub odpowiednią biblioteką.

\subsection{Kody błędów, wyjątki i wydajność}

W konwencjonalnych implementacjach C++ obsługa sukcesu na ścieżce,
na której żaden wyjątek nie wystąpił, może być bardzo tania, ale samo
zgłoszenie i obsłużenie wyjątku zwykle kosztuje więcej niż zwrócenie
prostego kodu błędu. Nie należy jednak zamieniać każdego \texttt{throw}
na kod zwrotny tylko na podstawie intuicji: warto mierzyć rzeczywistą
częstotliwość błędów, profilować aplikację i uwzględniać prostotę
kontraktów.

\section{Typowe błędy i dobre praktyki}
\label{sec:cpp11-guidelines}

\begin{enumerate}
\item \textbf{Nie rzucaj wskaźnikami do lokalnych obiektów.}
      Preferuj wyjątki będące obiektami, zwykle pochodnymi od
      \texttt{std::exception}.
\item \textbf{Przechwytuj wyjątki polimorficzne przez
      \texttt{const\&}.} Unikaj obcinania obiektu (\emph{slicing}).
\item \textbf{Do ponownego zgłoszenia używaj \texttt{throw;}.}
      Nie zastępuj go bez namysłu \texttt{throw e;}.
\item \textbf{Stosuj RAII.} Surowe wywołania \texttt{new}/\texttt{delete}
      rozdzielone operacjami rzucającymi wyjątki to częsta przyczyna
      wycieków pamięci.
\item \textbf{Nie wypuszczaj wyjątków z destruktorów.}
      Błędy finalizacji zgłaszaj jawnie w odpowiednio zaprojektowanym API.
\item \textbf{Nie oznaczaj funkcji \texttt{noexcept} bez analizy.}
      To obietnica z konsekwencją w postaci \texttt{std::terminate}.
\item \textbf{Nie ukrywaj błędów w pustym \texttt{catch (...)}.}
      Obsłuż wyjątek, zapisz informację diagnostyczną albo przekaż go wyżej.
\item \textbf{Określ gwarancję niepowodzenia.} Dla operacji mutujących
      dane rozważ silną gwarancję i strategię „przygotuj–zatwierdź”.
\item \textbf{Złap wyjątek w funkcji wątku.}
      Nie pozwól mu opuścić funkcji startowej \texttt{std::thread}.
      Do przekazania błędu użyj \texttt{std::exception\_ptr},
      \texttt{promise/future} lub API wyższego poziomu.
\item \textbf{Wybierz właściwy model błędu.} Brak zwykłego wyniku,
      przewidywalna porażka operacji i naruszenie kontraktu to różne
      sytuacje.
\end{enumerate}

\section{Przykład końcowy: transakcyjna aktualizacja danych}

Poniższy program ilustruje współpracę wyjątku, RAII i silnej gwarancji.
Załóżmy, że aktualizacja dodaje wartości do rejestru, ale w razie błędu
rejestr ma pozostać bez zmian.

\begin{lstlisting}
#include <iostream>
#include <stdexcept>
#include <utility>
#include <vector>

class Register {
    std::vector<int> values_;
public:
    Register() : values_{1, 2, 3} {}

    void add_checked(const std::vector<int>& numbers) {
        std::vector<int> proposed = values_;
        for (int x : numbers) {
            if (x < 0) {
                throw std::invalid_argument("negative input");
            }
            proposed.push_back(x);
        }
        values_.swap(proposed);
    }

    const std::vector<int>& values() const noexcept {
        return values_;
    }
};

int main() {
    Register r;
    try {
        r.add_checked({4, -5, 6});
    } catch (const std::invalid_argument& e) {
        std::cerr << e.what() << '\n';
    }
    // Still {1, 2, 3}; no partial update.
    for (int x : r.values()) std::cout << x << ' ';
}
\end{lstlisting}

Najpierw powstaje kopia danych. Sprawdzanie i alokowanie nowej pamięci
mogą zgłaszać wyjątki, lecz działają na kopii. Po pomyślnym wykonaniu
następuje zatwierdzenie przez \texttt{swap}. Jest to wzorzec, który
można stosować również w rozbudowanych klasach, przy świadomym
uwzględnieniu kosztów i własności ich składowych.

\section{Podsumowanie rozdziału}

W praktyce mechanizm wyjątków jest najbardziej użyteczny wtedy,
gdy łączy się go z dobrym modelem własności zasobów i z precyzyjnie
określonymi gwarancjami operacji. \texttt{throw}, \texttt{catch}
i \texttt{noexcept} to tylko składnia; o poprawności projektu
decydują RAII, niezmienniki obiektów, semantyka transakcyjna oraz
świadome oddzielanie błędów wyjątkowych od zwykłych wariantów wyniku.



\chapter{Wyjątki i bezpieczeństwo wyjątkowe w C++}
\label{chap:cpp-exceptions}

\section{Po co C++ ma wyjątki?}

Mechanizm wyjątków oddziela kod realizujący zwykły przebieg algorytmu
od kodu reagującego na niepowodzenia. Jeżeli funkcja nie może spełnić
swojego kontraktu, może zgłosić wyjątek poleceniem \texttt{throw}.
Wówczas wykonanie jest przekazywane do odpowiedniego bloku
\texttt{catch}; po drodze niszczone są obiekty automatyczne, których
czas życia dobiegł końca. To ostatnie zjawisko nazywamy
\emph{odwijaniem stosu} (\emph{stack unwinding}).

Najprostszy schemat wygląda tak: \texttt{throw} zgłasza problem,
a pasujący \texttt{catch} go odbiera. Kod po \texttt{throw} w tym samym
bloku nie zostanie wykonany:
\begin{lstlisting}
try {
    throw std::runtime_error("bad input");
} catch (const std::runtime_error& error) {
    std::cerr << error.what() << '\n';
}
\end{lstlisting}
W samodzielnym programie ten fragment wymaga nagłówków
\texttt{<iostream>} i \texttt{<stdexcept>}.

Z perspektywy projektowania API warto odróżniać:
\begin{itemize}
\item \emph{zwykły wynik} (w tym informację, że elementu nie znaleziono),
      dla którego często lepszy jest \texttt{std::optional};
\item \emph{błąd operacji, który wywołujący powinien świadomie obsłużyć},
      reprezentowany np. przez kod błędu lub, od C++23, \texttt{std::expected};
\item \emph{niepowodzenie uniemożliwiające dotrzymanie kontraktu},
      dla którego wyjątek bywa naturalnym sposobem propagacji.
\end{itemize}
Nie ma uniwersalnej reguły nakazującej używanie wyjątków wszędzie.
Znaczenie mają kontrakty funkcji, konwencje projektu, ograniczenia
środowiska oraz wymagania dotyczące opóźnień wykonania.

\section{Składnia \texttt{try}, \texttt{throw} i \texttt{catch}}

\begin{center}\small
\begin{tabularx}{\textwidth}{@{}p{0.32\textwidth}X@{}}
\toprule
Konstrukcja & Za co odpowiada \\
\midrule
\code{try} & Obejmuje blok, z którego wyjątki mogą zostać przechwycone przez następujące bloki \code{catch}. \\
\code{throw value} & Tworzy obiekt wyjątku i rozpoczyna poszukiwanie pasującej obsługi. \\
\texttt{catch(const E\& e)} & Przechwytuje wyjątek zgodnego typu bez kopiowania i bez obcinania części pochodnej. \\
\code{catch(...)} & Przechwytuje dowolny wyjątek, także nieznanego typu. \\
\code{throw;} & Wewnątrz obsługi ponownie zgłasza bieżący wyjątek, zachowując jego typ dynamiczny. \\
\code{std::exception::what()} & Zwraca tekstowy opis błędu; służy do diagnostyki, nie do sterowania logiką. \\
\bottomrule
\end{tabularx}
\end{center}
\label{sec:cpp11-basics}

\begin{lstlisting}
#include <iostream>
#include <stdexcept>

double divide(double a, double b) {
    if (b == 0.0)
        throw std::invalid_argument("zero denominator");
    return a / b;
}

int main() {
    try {
        std::cout << divide(10.0, 2.0) << '\n';
        std::cout << divide(10.0, 0.0) << '\n';
        std::cout << "not executed\n";
    } catch (const std::invalid_argument& e) {
        std::cerr << "invalid input: " << e.what() << '\n';
    }
}
\end{lstlisting}

Instrukcja \texttt{throw E\{...\}} tworzy obiekt wyjątku. Obsługa
szuka pasującej klauzuli \texttt{catch}; jeśli wyjątek nie zostanie
obsłużony, program kończy działanie poprzez \texttt{std::terminate}.
Nie należy zakładać, że nieobsłużony wyjątek zawsze spowoduje pełne
odwinięcie stosu: szczegóły zakończenia programu mogą zależeć od
implementacji.

\subsection{Przechwytywanie przez referencję i kolejność klauzul}

Przyjmujemy zwykle \texttt{const E\&}: pozwala to uniknąć niepotrzebnego
kopiowania i \emph{object slicing} przy przekazywaniu obiektu typu
pochodnego jako typu bazowego przez wartość.

\begin{lstlisting}
#include <exception>
#include <iostream>
#include <stdexcept>

void handle() {
    try {
        throw std::out_of_range("bad index");
    } catch (const std::out_of_range& e) {
        std::cout << "out_of_range: " << e.what() << '\n';
    } catch (const std::exception& e) {
        std::cout << "other std::exception: " << e.what() << '\n';
    } catch (...) {
        std::cout << "unknown exception\n";
    }
}
\end{lstlisting}

Klauzule są sprawdzane w kolejności występowania: najpierw umieszcza
się typy bardziej szczegółowe, a na końcu bazowe. Wyrażenie
\texttt{catch (...)} przechwytuje wyjątek dowolnego typu, ale nie
udostępnia bezpośrednio jego obiektu.

\paragraph{Ponowne zgłoszenie.}
\texttt{throw;} wewnątrz aktywnej obsługi ponownie zgłasza bieżący
wyjątek z zachowaniem jego dynamicznego typu. Natomiast
\texttt{throw e;} zgłasza nowe wyrażenie; w zależności od statycznego
typu \texttt{e} może prowadzić do kopiowania i obcięcia do klasy
bazowej.

\begin{lstlisting}
void log_and_rethrow() {
    try {
        throw std::runtime_error("disk error");
    } catch (const std::exception& e) {
        std::cerr << e.what() << '\n';
        throw;                  // keeps original exception object
        // throw e;              // may slice to std::exception
    }
}
\end{lstlisting}

\section{Odwijanie stosu i RAII}
\label{sec:cpp11-unwinding}

Najważniejsza reguła praktyczna brzmi: \emph{zwalnianie zasobu powinno
wynikać z destruktora obiektu zarządzającego zasobem}. Stanowi to
rozwinięcie idei RAII poznanej w rozdziale~\ref{chap:wskazniki-cpp}.

\begin{lstlisting}
#include <iostream>
#include <stdexcept>
#include <string>
#include <utility>

struct Trace {
    std::string name;
    explicit Trace(std::string n) : name(std::move(n)) {
        std::cout << "create " << name << '\n';
    }
    ~Trace() noexcept {
        std::cout << "destroy " << name << '\n';
    }
};

void third() {
    Trace c{"C"};
    throw std::runtime_error("failure");
}
void second() {
    Trace b{"B"};
    third();
}
void first() {
    Trace a{"A"};
    second();
}
// try { first(); } catch (const std::exception&) {}
// Destruction order: C, B, A.
\end{lstlisting}

Wyjątek powoduje zakończenie odpowiednich zakresów, a obiekty
z lokalnym automatycznym czasem przechowywania są niszczone w odwrotnej
kolejności ich skonstruowania. Jeżeli konstruktor zgłosi wyjątek,
destruktor \emph{całego, nieutworzonego obiektu} nie jest wywoływany,
ale niszczone są już w pełni skonstruowane podobiekty bazowe i składowe.

\subsection{Zarządzanie plikiem przy pomocy RAII}

\begin{lstlisting}
#include <cstdio>
#include <memory>
#include <stdexcept>

struct FileCloser {
    void operator()(std::FILE* file) const noexcept {
        if (file) std::fclose(file);
    }
};
using File = std::unique_ptr<std::FILE, FileCloser>;

File open_file(const char* name) {
    File file{std::fopen(name, "rb")};
    if (!file) throw std::runtime_error("cannot open file");
    return file;
}
\end{lstlisting}

Nawet gdy po otwarciu pliku dalsze operacje zgłoszą wyjątek,
\texttt{unique\_ptr} wykona deleter. Zauważmy, że
\texttt{std::fclose} może sygnalizować błąd przy zamykaniu; prosty
noexcept-deleter ignoruje ten kod błędu. Jeżeli informacja o błędzie
zamknięcia jest istotna, należy zapewnić osobną, jawną operację
finalizacji, której wynik można obsłużyć \emph{przed} destrukcją.

\section{Gwarancje bezpieczeństwa wyjątkowego}

\begin{center}\small
\begin{tabularx}{\textwidth}{@{}p{0.32\textwidth}X@{}}
\toprule
Gwarancja lub narzędzie & Skutek w razie wyjątku \\
\midrule
Brak wycieku & Zasoby pozostają poprawnie zwalniane, zwykle przez RAII. \\
Gwarancja podstawowa & Obiekty pozostają poprawne, choć ich stan może się zmienić. \\
Gwarancja silna & Operacja kończy się sukcesem albo pozostawia obserwowalny stan bez zmian. \\
Gwarancja braku wyjątków & Operacja nie pozwala wyjątkowi opuścić funkcji. \\
\code{noexcept} & Deklaruje gwarancję braku wyjątku; wyjątek opuszczający funkcję kończy program przez \code{std::terminate}. \\
\bottomrule
\end{tabularx}
\end{center}
\label{sec:cpp11-guarantees}

Bezpieczeństwo wyjątkowe nie oznacza tylko braku wycieków pamięci.
Załóżmy, że lista zawiera \texttt{[1, 2]} i dopisanie \texttt{3}
zgłosi wyjątek. Gwarancja podstawowa oznacza, że lista nadal nadaje
się do użycia, choć może mieć zmienioną zawartość. Gwarancja silna
oznacza, że nadal zawiera dokładnie \texttt{[1, 2]}. Gwarancja braku
wyjątków mówi tylko, że wyjątek nie opuści funkcji; jeśli operacja
ma dodatkowo zawsze osiągnąć cel, mówimy o gwarancji \emph{no-fail}.

Silna gwarancja nie oznacza automatycznie cofnięcia efektów
zewnętrznych, np. wysłanego pakietu czy wpisu w bazie danych.
Zakres transakcyjności trzeba sprecyzować w kontrakcie.

\subsection{Silna gwarancja dzięki operacji „przygotuj, a potem zatwierdź”}

\begin{lstlisting}
#include <utility>
#include <vector>

class Table {
    std::vector<int> data_;
public:
    explicit Table(std::vector<int> values)
        : data_(std::move(values)) {}

    void append_many(const std::vector<int>& more) {
        auto candidate = data_;          // may throw
        candidate.insert(candidate.end(),
                         more.begin(), more.end()); // may throw
        data_.swap(candidate);           // noexcept for std::vector<int>
    }

    const std::vector<int>& data() const noexcept { return data_; }
};
\end{lstlisting}

Jeżeli kopiowanie lub \texttt{insert} zawiedzie, \texttt{data\_}
pozostaje nietknięte. Dopiero bezwyjątkowe \texttt{swap} zatwierdza
przygotowaną wersję. Kosztem jest dodatkowa pamięć i kopiowanie;
silna gwarancja bywa więc świadomym wyborem projektowym, a nie darmową
optymalizacją.

\paragraph{Uwaga na aliasing.}
W przykładzie argument \texttt{more} może nawet odnosić się do
\texttt{data\_} obiektu \texttt{Table}: kopia \texttt{candidate}
powstaje przed modyfikacją oryginału. W kodzie operującym bezpośrednio
na tym samym kontenerze trzeba dodatkowo analizować ważność iteratorów
i referencji.

\subsection{Konstruktor i niezmienniki klasy}

Dobrze napisany konstruktor tworzy w pełni poprawny obiekt albo
zgłasza wyjątek. Bezpiecznie zarządzane składowe automatycznie
posprzątają po częściowo wykonanej inicjalizacji.

\begin{lstlisting}
#include <memory>
#include <stdexcept>

class Session {
    std::unique_ptr<int[]> buffer_;
    std::size_t size_;
public:
    explicit Session(std::size_t n)
        : buffer_(std::make_unique<int[]>(n)), size_(n) {
        if (n == 0) throw std::invalid_argument("empty session");
    }
    std::size_t size() const noexcept { return size_; }
};
\end{lstlisting}

Nie należy zastępować RAII ręcznie pisanymi \texttt{catch}, w których
sprząta się kolejne zasoby. Takie procedury są podatne na błędy,
zwłaszcza gdy konstruktor zdobywa wiele zasobów po kolei.

\section{\texttt{noexcept}: deklarowanie i sprawdzanie}
\label{sec:cpp11-noexcept}

\texttt{noexcept} w specyfikacji funkcji oznacza, że wyjątek nie może
opuścić jej ciała. Naruszenie tego wymagania kończy się wywołaniem
\texttt{std::terminate} (nie jest to „ciche zignorowanie” błędu).

\begin{lstlisting}
void close_handle() noexcept {
    // cleanup that must not let exceptions escape
}

void may_fail();
static_assert(noexcept(close_handle()));
static_assert(!noexcept(may_fail()));
\end{lstlisting}

W drugim użyciu \texttt{noexcept(expr)} jest \emph{operatorem}
zwracającym wartość logiczną w czasie kompilacji. Nie uruchamia on
samego wyrażenia.

\subsection{Warunkowe \texttt{noexcept}}

\begin{lstlisting}
#include <type_traits>
#include <utility>

template<class T>
void swap_values(T& a, T& b)
    noexcept(std::is_nothrow_move_constructible_v<T> &&
             std::is_nothrow_move_assignable_v<T>) {
    T tmp(std::move(a));
    a = std::move(b);
    b = std::move(tmp);
}
\end{lstlisting}

Jest to przykład \emph{warunkowej specyfikacji}. Nie oznacza,
że przedstawiony algorytm daje silną gwarancję, jeżeli przenoszenie
rzuca wyjątki; dla takich typów stan \texttt{a} lub \texttt{b} może
się zmienić przed niepowodzeniem. W kodzie generycznym trzeba
analizować nie tylko składnię, ale też gwarancje operacji typu \texttt{T}.

\subsection{Dlaczego \texttt{noexcept} jest ważne dla \texttt{vector}?}

Przy powiększaniu pojemności kontener może potrzebować przenieść
lub skopiować stare elementy do nowej pamięci. Dla typów kopiowalnych
konstrukcja kopiująca bywa bezpieczniejsza, jeśli konstruktor
przenoszący może zgłosić wyjątek.

\begin{lstlisting}
#include <type_traits>
#include <utility>

struct Item {
    Item() = default;
    Item(const Item&) {}              // copying available
    Item(Item&&) noexcept(false) {}   // move may throw
};

static_assert(std::is_copy_constructible_v<Item>);
static_assert(!std::is_nothrow_move_constructible_v<Item>);

// std::move_if_noexcept(x) usually gives const Item& here.
\end{lstlisting}

\texttt{std::move\_if\_noexcept} zwraca \texttt{T\&\&}, jeżeli
przenoszenie jest zadeklarowane jako bezwyjątkowe albo typ nie jest
kopiowalny; w przeciwnym razie zwraca \texttt{const T\&}.
Zasada ta pomaga kontenerom utrzymywać silną gwarancję w wielu,
ale nie we wszystkich przypadkach. Szczegóły są zależne od
konkretnej operacji i własności typu elementów.

\subsection{Destruktory i \texttt{std::terminate}}

Domyślnie destruktor jest niewyrzucający, o ile destruktory jego
składowych i klas bazowych nie wymagają innej specyfikacji.
Szczególnie podczas odwijania stosu wypuszczenie kolejnego wyjątku
z destruktora może spowodować \texttt{std::terminate}.

\begin{lstlisting}
struct SafeCleanup {
    ~SafeCleanup() noexcept {
        try {
            risky_cleanup();
        } catch (...) {
            // Log, record an error, or use another no-throw fallback.
            // Never let it escape this destructor.
        }
    }
    static void risky_cleanup();
};
\end{lstlisting}

Nie oznacza to, że każdą awarię należy ukryć. Jeśli operacja
finalizacji może się nie udać i użytkownik musi znać wynik, projektujmy
jawną metodę \texttt{close()} lub \texttt{commit()}, a destruktor
niech pełni rolę bezpiecznej ścieżki awaryjnej.

\section{Własne klasy wyjątków i hierarchia standardowa}
\label{sec:cpp11-types}

W nagłówkach \texttt{<stdexcept>}, \texttt{<exception>},
\texttt{<new>} i \texttt{<system\_error>} znajdziemy wiele gotowych
typów. W typowych zastosowaniach najczęściej spotkamy
\texttt{std::invalid\_argument}, \texttt{std::out\_of\_range},
\texttt{std::runtime\_error}, \texttt{std::bad\_alloc} oraz
\texttt{std::system\_error}.

\begin{lstlisting}
#include <cstddef>
#include <stdexcept>
#include <string>
#include <utility>

class ParseError : public std::runtime_error {
    std::size_t position_;
public:
    ParseError(std::size_t pos, std::string message)
        : std::runtime_error(std::move(message)),
          position_(pos) {}

    std::size_t position() const noexcept { return position_; }
};

int parse_digit(char c, std::size_t pos) {
    if (c < '0' || c > '9')
        throw ParseError(pos, "expected decimal digit");
    return c - '0';
}
\end{lstlisting}

Własna klasa może przechowywać dodatkowe informacje, np. identyfikator
operacji, pozycję w tekście lub kod błędu. Nie należy umieszczać tam
niepotrzebnie ogromnych danych. Metoda \texttt{what()} klasy
\texttt{std::exception} jest wirtualna i \texttt{noexcept}; jeśli
ją przesłaniamy, musimy zachować zgodny kontrakt i poprawny czas życia
zwracanego łańcucha.

\paragraph{Błędy systemowe.}
\texttt{std::system\_error} umożliwia przenoszenie kodu błędu
\texttt{std::error\_code} wraz z kategorią, np. błędu systemowego,
bez uzależniania całego API od przypadkowego tekstu \texttt{what()}.

\section{Wyjątki w funkcjach szablonowych i neutralność wyjątkowa}

Komponent generyczny powinien zwykle zachować niezmienniki i pozwolić
wyjątkom operacji wywołanych na typie \texttt{T} przejść do klienta,
zamiast pochłaniać je przez \texttt{catch (...)}.

\begin{lstlisting}
#include <algorithm>
#include <vector>

template<class T, class Compare>
void sort_elements(std::vector<T>& xs, Compare comp) {
    std::sort(xs.begin(), xs.end(), comp);
    // An exception from comp may propagate to the caller.
}
\end{lstlisting}

Należy jednak odróżnić \emph{neutralność wyjątkową} od silnej
gwarancji. Jeśli komparator zgłosi wyjątek w trakcie sortowania,
kontener pozostaje poprawnym obiektem, ale jego kolejność nie musi
odpowiadać stanowi sprzed wywołania.

\section{Przenoszenie wyjątków między funkcjami i wątkami}
\label{sec:cpp11-exception-ptr}

Ta i dwie kolejne sekcje pokazują zastosowania bardziej zaawansowane.
Do podstaw wystarczą \texttt{try}/\texttt{catch}, RAII oraz gwarancje
stanu po niepowodzeniu.

Typ \texttt{std::exception\_ptr} pozwala zachować wyjątek poza blokiem
\texttt{catch}, a później zgłosić go ponownie. Jest przydatny
szczególnie w programowaniu współbieżnym.

\begin{lstlisting}
#include <exception>
#include <iostream>
#include <stdexcept>
#include <string>
#include <thread>

std::exception_ptr error;
std::thread worker([&] {
    try {
        throw std::runtime_error("worker failed");
    } catch (...) {
        error = std::current_exception();
    }
});
worker.join();

if (error) {
    std::rethrow_exception(error);
}
\end{lstlisting}

Zmienna \texttt{error} nie jest tutaj współdzielona jednocześnie przez
dwa aktywne dostępy: zapis zachodzi w wątku roboczym, a odczyt
dopiero po \texttt{join()}, które zapewnia wymaganą synchronizację.
Nie wolno dopuścić, aby wyjątek opuścił funkcję startową zwykłego
\texttt{std::thread}: typowym skutkiem jest \texttt{std::terminate}.

\subsection{\texttt{std::future} przechowuje wyjątek}

Wspomniane w rozdziale o współbieżności \texttt{std::async} i
\texttt{std::future} ułatwiają propagowanie niepowodzenia do
odbiorcy wyniku.

\begin{lstlisting}
#include <future>
#include <stdexcept>

auto fut = std::async(std::launch::async, []() -> int {
    throw std::runtime_error("calculation failed");
});
try {
    (void)fut.get();  // rethrows the stored exception
} catch (const std::runtime_error&) {
    // report the failure in the receiving thread
}
\end{lstlisting}

Przez \texttt{get()} wynik (lub wyjątek) trafia do wątku odbiorcy.
W \texttt{std::future} wynik można pobrać tylko raz; inne własności
tego mechanizmu omówiliśmy w rozdziale o asynchroniczności.

\section{Wyjątki zagnieżdżone i dodatkowy kontekst}
\label{sec:cpp11-nested}

Gdy funkcja wyższego poziomu rozumie więcej niż warstwa niższa, może
wzbogacić informację o błędzie i zachować oryginalną przyczynę.

\begin{lstlisting}
#include <exception>
#include <stdexcept>

void load_config() {
    try {
        throw std::runtime_error("invalid token");
    } catch (...) {
        std::throw_with_nested(
            std::runtime_error("cannot load configuration"));
    }
}

void print_nested(const std::exception& e, int depth = 0) {
    std::cerr << std::string(depth * 2, ' ') << e.what() << '\n';
    try {
        std::rethrow_if_nested(e);
    } catch (const std::exception& inner) {
        print_nested(inner, depth + 1);
    } catch (...) {
        std::cerr << "nested non-standard exception\n";
    }
}
\end{lstlisting}

Takie opakowywanie wyjątków nie jest tym
samym co zgłoszenie \texttt{throw e;}: nowy wyjątek przechowuje
kontekst, a przyczynę można odzyskać przez \texttt{rethrow\_if\_nested}.

\section{\texttt{std::uncaught\_exceptions} i strażnicy zakresu}

Funkcja \texttt{std::uncaught\_exceptions()} (C++17) zwraca liczbę
aktualnie nieobsługiwanych obiektów wyjątków w bieżącym wątku.
Można jej użyć do wykrycia, czy wyjście z zakresu następuje wskutek
odwijania stosu.

\begin{lstlisting}
#include <exception>

class FailureGuard {
    int initial_ = std::uncaught_exceptions();
    int& failure_count_;
public:
    explicit FailureGuard(int& count) noexcept : failure_count_(count) {}
    FailureGuard(const FailureGuard&) = delete;
    FailureGuard& operator=(const FailureGuard&) = delete;

    ~FailureGuard() noexcept {
        if (std::uncaught_exceptions() > initial_) {
            ++failure_count_; // no-throw bookkeeping
        }
    }
};
\end{lstlisting}

Porównanie z \emph{początkową liczbą}, a nie tylko z zerem, jest
ważne przy zagnieżdżonych wyjątkach. Strażnik nie daje jednak
pełnego, transakcyjnego mechanizmu cofania zmian. W projektach
korzystających z C++20 można napisać własny scope guard, ale trzeba
uważnie określić, czy przechowuje funkcję przez kopię, przez
przeniesienie czy przez referencję i czy jej wywołanie może rzucać.

\section{Wyjątki a alternatywne modele obsługi błędów}
\label{sec:cpp11-alternatives}

\subsection{\texttt{std::optional}: brak wyniku nie musi być błędem}

\begin{lstlisting}
#include <optional>
#include <string_view>

std::optional<std::size_t> find_colon(std::string_view s) {
    const auto i = s.find(':');
    if (i == std::string_view::npos) return std::nullopt;
    return i;
}
\end{lstlisting}

Brak dwukropka może być normalnym scenariuszem. Zgłaszanie wyjątku
przy każdym nieudanym wyszukiwaniu bywa niepotrzebne.

\subsection{\texttt{std::expected<T,E>} — C++23, nie C++20}

\texttt{std::expected} reprezentuje \emph{albo wartość typu T, albo
błąd typu E}. Został ustandaryzowany w C++23 i nie należy zakładać,
że nagłówek \texttt{<expected>} jest dostępny w kompilacji C++20.

\begin{lstlisting}
// This example requires C++23 and <expected>.
#include <expected>
#include <string>
#include <string_view>

std::expected<int, std::string>
parse_positive(std::string_view text) {
    if (text == "1") return 1;
    if (text == "2") return 2;
    return std::unexpected("expected 1 or 2");
}

void use_expected() {
    auto result = parse_positive("3");
    if (!result) {
        const std::string& error = result.error();
        (void)error;
    }
}
\end{lstlisting}

Dzięki temu błędy oczekiwane przez API są jawne w sygnaturze.
\texttt{std::expected} nie eliminuje wszystkich wyjątków: konstrukcja
zawartego obiektu czy inne operacje nadal mogą się nie udać.
W C++20 podobny kontrakt można zrealizować własnym typem
\texttt{Result}, \texttt{std::variant<T,E>} lub odpowiednią biblioteką.

\subsection{Kody błędów, wyjątki i wydajność}

W konwencjonalnych implementacjach C++ obsługa sukcesu na ścieżce,
na której żaden wyjątek nie wystąpił, może być bardzo tania, ale samo
zgłoszenie i obsłużenie wyjątku zwykle kosztuje więcej niż zwrócenie
prostego kodu błędu. Nie należy jednak zamieniać każdego \texttt{throw}
na kod zwrotny tylko na podstawie intuicji: warto mierzyć rzeczywistą
częstotliwość błędów, profilować aplikację i uwzględniać prostotę
kontraktów.

\section{Typowe błędy i dobre praktyki}
\label{sec:cpp11-guidelines}

\begin{enumerate}
\item \textbf{Nie rzucaj wskaźnikami do lokalnych obiektów.}
      Preferuj wyjątki będące obiektami, zwykle pochodnymi od
      \texttt{std::exception}.
\item \textbf{Przechwytuj wyjątki polimorficzne przez
      \texttt{const\&}.} Unikaj obcinania obiektu (\emph{slicing}).
\item \textbf{Do ponownego zgłoszenia używaj \texttt{throw;}.}
      Nie zastępuj go bez namysłu \texttt{throw e;}.
\item \textbf{Stosuj RAII.} Surowe wywołania \texttt{new}/\texttt{delete}
      rozdzielone operacjami rzucającymi wyjątki to częsta przyczyna
      wycieków pamięci.
\item \textbf{Nie wypuszczaj wyjątków z destruktorów.}
      Błędy finalizacji zgłaszaj jawnie w odpowiednio zaprojektowanym API.
\item \textbf{Nie oznaczaj funkcji \texttt{noexcept} bez analizy.}
      To obietnica z konsekwencją w postaci \texttt{std::terminate}.
\item \textbf{Nie ukrywaj błędów w pustym \texttt{catch (...)}.}
      Obsłuż wyjątek, zapisz informację diagnostyczną albo przekaż go wyżej.
\item \textbf{Określ gwarancję niepowodzenia.} Dla operacji mutujących
      dane rozważ silną gwarancję i strategię „przygotuj–zatwierdź”.
\item \textbf{Złap wyjątek w funkcji wątku.}
      Nie pozwól mu opuścić funkcji startowej \texttt{std::thread}.
      Do przekazania błędu użyj \texttt{std::exception\_ptr},
      \texttt{promise/future} lub API wyższego poziomu.
\item \textbf{Wybierz właściwy model błędu.} Brak zwykłego wyniku,
      przewidywalna porażka operacji i naruszenie kontraktu to różne
      sytuacje.
\end{enumerate}

\section{Przykład końcowy: transakcyjna aktualizacja danych}

Poniższy program ilustruje współpracę wyjątku, RAII i silnej gwarancji.
Załóżmy, że aktualizacja dodaje wartości do rejestru, ale w razie błędu
rejestr ma pozostać bez zmian.

\begin{lstlisting}
#include <iostream>
#include <stdexcept>
#include <utility>
#include <vector>

class Register {
    std::vector<int> values_;
public:
    Register() : values_{1, 2, 3} {}

    void add_checked(const std::vector<int>& numbers) {
        std::vector<int> proposed = values_;
        for (int x : numbers) {
            if (x < 0) {
                throw std::invalid_argument("negative input");
            }
            proposed.push_back(x);
        }
        values_.swap(proposed);
    }

    const std::vector<int>& values() const noexcept {
        return values_;
    }
};

int main() {
    Register r;
    try {
        r.add_checked({4, -5, 6});
    } catch (const std::invalid_argument& e) {
        std::cerr << e.what() << '\n';
    }
    // Still {1, 2, 3}; no partial update.
    for (int x : r.values()) std::cout << x << ' ';
}
\end{lstlisting}

Najpierw powstaje kopia danych. Sprawdzanie i alokowanie nowej pamięci
mogą zgłaszać wyjątki, lecz działają na kopii. Po pomyślnym wykonaniu
następuje zatwierdzenie przez \texttt{swap}. Jest to wzorzec, który
można stosować również w rozbudowanych klasach, przy świadomym
uwzględnieniu kosztów i własności ich składowych.

\section{Podsumowanie rozdziału}

W praktyce mechanizm wyjątków jest najbardziej użyteczny wtedy,
gdy łączy się go z dobrym modelem własności zasobów i z precyzyjnie
określonymi gwarancjami operacji. \texttt{throw}, \texttt{catch}
i \texttt{noexcept} to tylko składnia; o poprawności projektu
decydują RAII, niezmienniki obiektów, semantyka transakcyjna oraz
świadome oddzielanie błędów wyjątkowych od zwykłych wariantów wyniku.

